\documentclass[annotation,times]{itmo-student-thesis}

%% Опции пакета:
%% - specification - если есть, генерируется задание (здесь отключено)
%% - annotation    - если есть, генерируется аннотация
%% - times         - шрифт Times New Roman, сборка через xelatex

\usepackage{icomma}
\usepackage{tabularx}
\usepackage{longtable}
\usepackage{booktabs}

\addbibresource{references.bib}

\begin{document}

\studygroup{M4239}
\title{Гибридное колоночно-разреженное хранение полуструктурированных данных в СУБД Otterbrix}
\author{Малько Егор}{Малько Е.}
\supervisor{}{}{}{}
\publishyear{2026}
\startdate{01}{сентября}{2024}
\finishdate{31}{мая}{2026}
\defencedate{15}{июня}{2026}

\secretary{}

%%% Цель исследования (требуется опцией annotation)
\researchaim{Разработка эффективного механизма хранения
полуструктурированных данных в рамках колоночного движка
Otterbrix, сочетающего производительность колоночного доступа к
часто встречающимся полям с компактным представлением редко
встречающихся полей.}

%%% Задачи, решаемые в ВКР (требуется опцией annotation)
\researchtargets{\begin{enumerate}
    \item проанализировать существующие подходы к хранению
    полуструктурированных данных в СУБД;
    \item спроектировать гибридное представление, сочетающее
    колоночное хранение плотных полей и компактное хранение
    редких полей с автоматической миграцией между режимами;
    \item реализовать предложенное представление в кодовой базе
    Otterbrix;
    \item обеспечить корректное поведение системы при параллельных
    операциях INSERT и SELECT над таблицами с вычисляемой схемой;
    \item выполнить оптимизацию физических операторов для
    достижения производительности, сопоставимой с существующими
    промышленными решениями;
    \item провести сравнительное тестирование на стандартном
    бенчмарке JSONBench с ClickHouse и PostgreSQL.
\end{enumerate}}

%%% Использование современных пакетов программ и технологий
\addadvancedsoftware{Язык {C++20} с корутинами, полиморфными
аллокаторами и концептами}{Глава~\ref{ch:impl}}
\addadvancedsoftware{Библиотека Boost.JSON для разбора
JSON-документов}{Глава~\ref{ch:impl}}
\addadvancedsoftware{Актёрная модель actor-zeta для
асинхронного исполнения запросов}{Глава~\ref{ch:impl}}

%%% Краткая характеристика полученных результатов (для аннотации)
\researchsummary{Спроектирован и реализован механизм гетерогенного
хранения полуструктурированных данных в колоночном движке Otterbrix:
часто встречающиеся поля JSON хранятся как полноценные физические
колонки, редкие~--- во вспомогательных двухколоночных таблицах вида
\texttt{(\_id, value)}; при пересечении заранее зафиксированного
порога заполненности поле автоматически мигрирует в основную таблицу.
Чтение разреженных колонок реализовано через реконструкцию на
уровне диспетчера, семантически эквивалентную LEFT~JOIN. Выполнены
оптимизации физических операторов: проекция колонок через
оптимизатор, типизированные предикаты фильтрации, специализированные
пути GROUP BY, исправление zone~maps. На стандартном бенчмарке
JSONBench система показывает производительность, сопоставимую с
ClickHouse 24.10, и существенно превосходит PostgreSQL с JSONB.}

%%% Гранты
\researchfunding{Работа выполнялась в рамках участия в проекте
Otterbrix с открытым исходным кодом без внешнего финансирования.}

%%% Публикации
\researchpublications{По теме работы публикаций не имеется.}

%% Генерируем титульный лист и аннотацию
\maketitle{Магистр}

%% Оглавление
\tableofcontents

%% Введение
\startprefacepage
%% Введение

Полуструктурированные данные в формате JSON прочно заняли место
основного способа обмена данными между веб-сервисами, форматом
логирования событий, телеметрии устройств и хранилищем данных
социальных сетей. Одновременно традиционные реляционные СУБД
оказались плохо приспособлены к такого рода нагрузкам: хранение
JSON в виде неделимого значения требует полного считывания
документа с диска при обращении к единственному полю, а
классическое колоночное хранение нацелено на заранее известную
фиксированную схему.

Современные промышленные СУБД решают эту задачу по-разному.
Документоориентированные системы (MongoDB как исторический
законодатель этого направления, Couchbase) оптимизируют работу с
произвольной структурой за счёт отказа от реляционного интерфейса
и колоночного хранения. Колоночные аналитические СУБД (ClickHouse,
Snowflake, DuckDB) вводят специализированные типы данных,
сочетающие заранее объявленные колонки и общую область хранения
редких полей. Каждый из подходов сопряжён с компромиссами:
документные системы не обеспечивают производительности
аналитических запросов, а специализированные колоночные форматы
ограничены в выборе операций над редкими полями.

Настоящая работа посвящена разработке альтернативного механизма
хранения полуструктурированных данных в рамках колоночной СУБД
Otterbrix~--- проекта с открытым исходным кодом, развивающегося как
встраиваемый аналитический движок с поддержкой SQL. Отличительная
особенность предлагаемого решения заключается в сочетании колоночного
хранения частых полей и размещения редких полей в отдельных
вспомогательных двухколоночных таблицах, связанных с основной через
логический идентификатор. Такая архитектура даёт возможность
переиспользовать всю механику колоночного движка (векторизация,
zone~maps, индексы, MVCC, WAL) для редких полей без необходимости
поддерживать отдельную подсистему хранения.

В рамках работы были решены следующие задачи. Во-первых, выполнен
обзор существующих подходов к хранению JSON в СУБД: BSON в
MongoDB как канонический пример документной модели, модель
Entity-Attribute-Value и её варианты, колоночное развёртывание
путей, новый тип \texttt{JSON} в ClickHouse. Во-вторых, предложена
модель гетерогенного хранения, в которой статус поля (закреплённое
в основной таблице или вынесенное в отдельную side-таблицу)
фиксируется при \texttt{CREATE TABLE}, а множество полей таблицы
динамически расширяется по мере появления новых ключей в
поступающих JSON-документах. В-третьих, модель реализована в
кодовой базе Otterbrix: введён класс \texttt{computed\_schema},
инкапсулирующий динамическую схему, добавлены четыре физических
оператора (\texttt{scan\_computing\_table} для гибридного
сканирования и три варианта DML \texttt{*\_computing} для
multi-collection-операций под одной транзакцией), pre-flight в
диспетчере для DDL side-таблиц, поддержка SQL-конструкций
\texttt{CREATE TABLE db.t ();},
\texttt{WITH (pinned\_columns=\ldots)} и
\texttt{INSERT \ldots\ VALUES (json('\ldots'))}. В-четвёртых,
выполнены оптимизации физических операторов: проекция колонок на
уровне оптимизатора, типизированные предикаты фильтрации с
пакетным копированием подошедших строк, специализированные пути
агрегации, исправление ошибки в реализации zone~maps. В-пятых,
проведено сравнительное тестирование на стандартном бенчмарке
JSONBench: на трёх конфигурациях самой Otterbrix (исходная
документная, новая колоночно-разреженная, новая
колоночно-разреженная с оптимизациями) и в сравнении с
встраиваемым колоночным движком DuckDB.

Результаты показали, что архитектурный переход от документной
модели к колоночно-разреженной даёт ускорение в 3--14~раз;
оптимизации физических операторов добавляют ещё 4--23~раза.
Финальная конфигурация Otterbrix на запросах JSONBench
сопоставима с DuckDB на аналитических запросах по плотным полям
(отставание единицы процентов) и обходит DuckDB в 1.3--3.6~раз на
запросах с полным сканированием и обращением к разреженным полям.
Объём новой кодовой базы измеряется тысячами строк~--- на порядок
меньше того, что потребовался бы для разработки специализированной
подсистемы хранения с нуля.

Настоящая работа организована следующим образом. В главе~1
изложена постановка задачи, описана исходная архитектура Otterbrix
и мотивация выбранного направления. Глава~2 содержит обзор
существующих подходов. В главе~3 описывается спроектированная
модель хранения. Глава~4 посвящена деталям реализации. В главе~5
рассматриваются оптимизации физических операторов. Глава~6
приводит результаты сравнительного тестирования с DuckDB и с
тремя конфигурациями самой Otterbrix.


%% Основные главы
\chapter{Постановка задачи и исходная архитектура Otterbrix}

\section{Общая характеристика проекта}

Otterbrix~--- встраиваемый вычислительный фреймворк с открытым
исходным кодом, предназначенный для обработки полуструктурированных
данных в сочетании с классической реляционной моделью. Проект
развивается как альтернатива специализированным документным СУБД
(MongoDB, Couchbase) и аналитическим колоночным системам
(ClickHouse, DuckDB), предлагая унифицированную модель, в которой
SQL-запросы могут одновременно обращаться к структурированным
таблицам и к полям JSON-документов в рамках одного исполнителя.

Архитектурно Otterbrix построен на следующих принципах современных
СУБД:
\begin{itemize}
    \item \textbf{Акторная модель параллелизма} (реализована поверх
    библиотеки actor-zeta)~--- все ключевые компоненты (дисковый
    менеджер, WAL, диспетчер запросов, менеджер индексов)
    оформлены как независимые актёры, обменивающиеся сообщениями.
    Это даёт предсказуемую изоляцию состояния и позволяет
    масштабировать систему на большое число ядер без глобальных
    блокировок.
    \item \textbf{Корутины C++20} (co\_await~/~co\_return)~---
    используются для организации асинхронного исполнения планов
    запросов. Каждый план может приостанавливаться на ожидании
    диска или других актёров, не блокируя нить пула.
    \item \textbf{Полиморфные аллокаторы} (\texttt{std::pmr})~---
    вся работа с памятью идёт через конкретный
    \texttt{memory\_resource}, что позволяет локализовать
    аллокации запросов и контролировать фрагментацию.
    \item \textbf{Собственный SQL-парсер на базе
    libpg\_query}~--- грамматика PostgreSQL обеспечивает
    совместимость с широким спектром существующих запросов и
    инструментов.
    \item \textbf{Колоночное хранилище}~--- таблицы организованы
    как последовательность row~groups, каждый из которых содержит
    независимые колоночные сегменты. Это стандартный приём для
    ускорения аналитических запросов.
\end{itemize}

На момент начала работы проект находился в активной фазе
рефакторинга: команда завершала переход с гибридной модели
хранения (отдельные подсистемы для документов и таблиц) на единое
колоночное хранилище. Этот переход, оформленный одним крупным
коммитом (\texttt{a34d3823}, <<Remove components/document and
migrate to columnar data\_chunk storage>>, февраль 2026), удалил
более 18~000 строк кода и послужил отправной точкой для настоящей
работы.

\section{Двойственность представления: документы и таблицы}

До указанного рефакторинга Otterbrix поддерживал два
принципиально различных типа коллекций.

\textbf{Документные коллекции} использовали отдельный компонент
\texttt{components/document/}, реализовавший собственный вариант
иерархического JSON-подобного значения. Каждый документ
представлял собой дерево (объекты, массивы, скалярные листья),
сериализованное в формате, близком к BSON. Доступ к полям
происходил по JSON-указателю (\texttt{/commit/record/text}),
каждый документ хранился как неделимое значение.

\textbf{Табличные коллекции} использовали колоночное хранилище
\texttt{components/table/} с фиксированной схемой (набор колонок
фиксированных типов). Доступ к ячейкам происходил по паре
(row\_id, column\_id); каждая колонка хранилась в отдельном
сегменте памяти или диска.

Обе реализации имели собственную подсистему хранения, собственные
операторы исполнения (match, group, sort, insert, update, delete),
собственные адаптеры диска, собственный WAL-формат. Класс
\texttt{document\_t} (\texttt{components/document/document.hpp},
удалён в \texttt{a34d3823}) насчитывал около 1100~строк кода и
реализовывал полнофункциональный JSON-документ с поддержкой
мутаций через \texttt{set(json\_pointer, value)},
\texttt{remove(json\_pointer)}, \texttt{set\_array},
\texttt{set\_dict}, а также операций сравнения и сериализации.

Основные недостатки документной реализации:
\begin{itemize}
    \item \textbf{Неэффективный ввод-вывод при аналитических
    запросах.} Для ответа на запрос вида \texttt{SELECT COUNT(*)
    FROM t WHERE commit.operation = 'create'} система вынуждена
    была читать с диска полный документ целиком, хотя для ответа
    требовалось единственное поле.
    \item \textbf{Низкая локальность кэша L1/L2.} Каждое
    обращение по JSON-указателю представляло собой обход дерева
    узлов, расположенных в куче и связанных указателями.
    \item \textbf{Невозможность применения векторных
    оптимизаций.} Векторизованные операторы, повсеместно
    используемые в современных колоночных СУБД (пакетное
    сравнение, SIMD), требуют компактного, однотипного
    расположения значений в памяти.
    \item \textbf{Удвоенная кодовая база.} Наличие двух
    параллельных реализаций операторов приводило к постоянной
    рассинхронизации: оптимизации одного пути не переносились в
    другой; баги всплывали в двух местах независимо.
\end{itemize}

\section{Миграция на единое колоночное хранилище}

Коммит \texttt{a34d3823} завершил миграцию: компонент
\texttt{components/document/} был удалён полностью, все коллекции
стали хранить данные в формате \texttt{data\_chunk\_t}~---
структуры из \texttt{components/vector/}, представляющей собой
вектор колонок с общим числом строк. Каждая колонка хранится в
виде сегментов \texttt{vector\_t}, поддерживающих несколько
представлений (flat, constant, dictionary) и плоский массив
значений одного типа.

После миграции таблицы получили два подвида:
\begin{itemize}
    \item \textbf{Таблицы с фиксированной схемой.} Колонки
    задаются при \texttt{CREATE TABLE db.t (id INT, name TEXT);}
    и остаются неизменными в течение жизни таблицы.
    \item \textbf{Таблицы с вычисляемой (computed) схемой.}
    Создаются без предварительного перечисления колонок~---
    \texttt{CREATE TABLE db.t ();}. Каждая операция
    \textsc{insert} способна добавить новые поля; схема
    расширяется по мере поступления данных. Именно этот режим
    заменил бывшие <<документные>> коллекции и стал основным
    предметом настоящей работы.
\end{itemize}

Миграция закрыла проблему удвоенной кодовой базы, но оставила
открытым ключевой вопрос: \textbf{как именно хранить колонки,
заполненные лишь в небольшой доле строк}. Прямолинейный вариант~---
создавать физическую колонку на каждый путь JSON и заполнять её
NULL для отсутствующих значений~--- порождает две новые проблемы:
\begin{enumerate}
    \item \textbf{Катастрофический рост пространства.} Реальные
    полуструктурированные данные (в том числе стандартный датасет
    JSONBench, сформированный из публичного потока сети Bluesky)
    содержат десятки тысяч различных JSON-путей, большинство из
    которых встречаются в менее чем 0.1\% документов.
    \item \textbf{Замедление записи.} Каждое добавление новой
    колонки в существующую таблицу требует выделения буфера для
    всех $N$ ранее вставленных строк и заполнения его значениями
    по умолчанию.
\end{enumerate}

\section{Формулировка задачи}

Цель выпускной квалификационной работы сформулирована следующим
образом: спроектировать и реализовать эффективный механизм
хранения полуструктурированных данных в рамках колоночного
движка Otterbrix, сочетающий производительность колоночного
доступа к часто встречающимся полям с компактным представлением
редко встречающихся полей, без необходимости поддерживать
отдельную подсистему для JSON-документов.

Для достижения поставленной цели решаются следующие задачи:
\begin{enumerate}
    \item Проанализировать существующие подходы к хранению
    полуструктурированных данных (BSON в PostgreSQL, EAV,
    flatten-представление путей, тип \texttt{JSON} в ClickHouse).
    \item Спроектировать гибридное представление, сочетающее
    колоночное хранение плотных полей и компактное хранение
    редких полей, с автоматической миграцией между режимами.
    \item Реализовать предложенное представление в кодовой базе
    Otterbrix, встраиваясь в существующие механизмы SQL-парсера,
    диспетчера и подсистемы хранения.
    \item Обеспечить корректное поведение системы при
    параллельных операциях \textsc{insert} и \textsc{select} над
    таблицами с вычисляемой схемой.
    \item Выполнить оптимизацию физических операторов (проекция
    колонок, типизированные предикаты, векторизация) для
    достижения производительности, сопоставимой с существующими
    промышленными решениями.
    \item Провести сравнительное тестирование на бенчмарке
    JSONBench с ClickHouse и PostgreSQL в качестве референсных
    систем.
\end{enumerate}

\section{Научная новизна и практическая значимость}

Основная новизна работы заключается в предложении механизма
\textbf{гетерогенного хранения колонок}: в рамках одной логической
таблицы часть колонок хранится в основной таблице (колоночное
хранение), часть~--- в отдельных вспомогательных таблицах формата
\texttt{(\_id, value)} (разреженное хранение), с автоматической
миграцией полей между режимами по мере накопления данных. В отличие
от подхода ClickHouse, где редкие поля попадают во внутренний
<<shared~data>> формат с ограниченной функциональностью,
предложенное решение сохраняет для разреженных колонок полный
набор операций (произвольные предикаты, индексы, агрегации) за счёт
реализации чтения через LEFT~JOIN на уровне исполнителя.

Практическая значимость состоит в получении системы, которая
обходит известные ограничения наивного колоночного хранения JSON
(неограниченный рост схемы, раздувание нулевыми значениями), не
требует поддержки отдельной подсистемы для документов, демонстрирует
производительность, сопоставимую с ClickHouse и существенно
превосходящую PostgreSQL с JSONB, и реализована на основе
существующих в Otterbrix механизмов колоночного хранения, что
позволяет переиспользовать все накопленные оптимизации
(векторизация, zone~maps, projection~pushdown) без дополнительных
усилий.

\chapterconclusion

В настоящей главе описана исходная архитектура проекта Otterbrix,
мотивация отказа от документной подсистемы хранения в пользу
единой колоночной архитектуры, и сформулированы задачи, решаемые
в настоящей выпускной квалификационной работе. Последующие главы
посвящены обзору существующих подходов, проектированию и реализации
предлагаемого решения.

\chapter{Исследование существующих подходов}

\startrelatedwork

Настоящая глава содержит обзор существующих способов хранения
полуструктурированных данных в СУБД. Рассматриваются четыре
принципиально различных подхода: хранение JSON как неделимого
значения, модель Entity-Attribute-Value, колоночное развёртывание
путей, и гибридный тип \texttt{JSON} в ClickHouse.

Каждый из подходов анализируется по трём характеристикам:
эффективность ввода-вывода при запросе к отдельному полю,
локальность кэша при обработке большого числа строк, и стоимость
эволюции схемы при появлении нового поля.

\section{Хранение JSON как неделимого значения}

Исторически первый и до сих пор наиболее распространённый способ
хранения JSON в СУБД~--- рассматривать документ целиком как
значение одной ячейки особого типа. Документ сериализуется в
текстовом виде (JSON) или в бинарном (BSON в MongoDB, JSONB в
PostgreSQL).

Текстовый JSON (MySQL до 5.7, ранние версии Oracle) хранит JSON как
тип CLOB без внутреннего индексирования: при каждом обращении к полю
происходит полный разбор строки парсером.

Каноническим представителем документной модели на сегодняшний день
является MongoDB~\cite{mongodb-bson}, разработавшая бинарный формат
BSON и фактически закрепившая в индустрии саму концепцию
<<документоориентированной СУБД>>. В BSON документ предварён
таблицей смещений внутренних элементов, каждое скалярное значение
хранится в нативном бинарном формате, типы значений
самоописательны. Близкое по идее представление JSONB в PostgreSQL
дополнительно упорядочивает ключи объектов для ускорения бинарного
поиска. Оба подхода позволяют обращаться к отдельному полю за время
$O(\log n)$ без полного разбора, однако два фундаментальных
недостатка сохраняются:

\begin{enumerate}
    \item \textbf{Полное чтение документа с диска.} Страница диска в
    СУБД типично составляет 4--16~КБ, а один JSON-документ нередко
    имеет размер от единиц до десятков килобайт. Для запроса
    \texttt{SELECT doc.user.id FROM t WHERE doc.kind = 'commit'}
    система должна считать с диска каждый документ целиком.
    \item \textbf{Отсутствие колоночного доступа.} Несколько
    последовательных строк таблицы содержат свои документы,
    разнесённые по страницам; чтение одного и того же поля для
    миллиона строк приводит к миллиону разрозненных обращений к
    памяти.
\end{enumerate}

Эмпирически на аналитических нагрузках над JSON-данными
документная модель и JSONB демонстрируют производительность на
один-два порядка хуже колоночных СУБД на том же объёме данных.
Индексы (составные, функциональные, в случае MongoDB~--- в том
числе wildcard) частично ускоряют выборку по конкретному пути, но
не решают фундаментальной проблемы ввода-вывода при последующем
чтении других полей и не ускоряют агрегатные запросы,
затрагивающие всю таблицу. Именно эти ограничения сделали
документную модель плохо приспособленной к OLAP-нагрузкам и
послужили мотивацией для появления специализированных колоночных
форматов хранения JSON.

\section{Модель Entity-Attribute-Value}

EAV~--- классическая реляционная техника моделирования данных с
переменной структурой, описанная в литературе ещё в 1970--1980-е
годы в контексте медицинских и биологических баз
данных~\cite{stead1983,nadkarni2007}. Идея состоит в декомпозиции
гетерогенных записей в универсальную тройку:
\begin{verbatim}
CREATE TABLE attr (
  entity_id BIGINT,
  attr_name TEXT,
  attr_value TEXT
);
\end{verbatim}

Документ \texttt{\{"a":~1, "b":~"x"\}} превращается в две строки:
\texttt{(1, 'a', '1')} и \texttt{(1, 'b', 'x')}. Запрос
\texttt{SELECT a FROM t} принимает вид
\texttt{SELECT attr\_value FROM attr WHERE attr\_name~=~'a'}.

Сильные стороны: схема полностью динамическая, отсутствие
NULL-значений, простая модель мутаций. Слабые стороны, многократно
описанные в литературе~\cite{dinu2007}:

\begin{itemize}
    \item \textbf{Потеря типизации.} Все значения приводятся к
    строке. Агрегация числовых полей требует преобразования типов
    на каждой строке.
    \item \textbf{Деградация запросов с несколькими полями.} Запрос
    \texttt{SELECT a, b, c FROM t} становится трёхкратным
    self-join либо группировкой с условными агрегатами. Сложность
    превышает колоночный вариант на порядок.
    \item \textbf{Иерархия.} Для вложенных JSON-документов
    возникает выбор между двумя не-эквивалентными решениями:
    через \texttt{parent\_id} (рекурсивный обход) и через
    разворот путей (flattened~paths, \cite{beyer2005,florescu1999}).
\end{itemize}

Обе вариации EAV страдают от общего недостатка: колоночный доступ к
одному полю невозможен без предикатного сканирования по
\texttt{attr\_name}. Как следствие, аналитические нагрузки
деградируют многократно. Это явление формулируется в работах по
декомпозиции хранилищ как <<vertical~fragmentation>> проблема
групповой обработки атрибутов~\cite{copeland1985}.

\section{Колонка на каждый путь JSON}

Альтернативный подход, к которому пришли современные аналитические
движки (Snowflake VARIANT~\cite{snowflake-variant}, ClickHouse JSON
в ранних версиях, ряд коммерческих систем)~--- разворачивать
JSON-документ в колоночное представление, где каждый уникальный
путь становится отдельной физической колонкой.

Документ \texttt{\{"a":~1, "b":~"x", "c":\{"d":~2\}\}} при вставке
порождает три колонки: \texttt{a}~(INT), \texttt{b}~(STRING),
\texttt{c.d}~(INT). Отсутствующие в конкретном документе поля
записываются как NULL.

Сильные стороны: колоночный доступ в полной мере, полная типизация,
применимость всех стандартных оптимизаций (сжатие, zone~maps,
векторизация).

Слабые стороны:

\begin{itemize}
    \item \textbf{Рост числа физических колонок.} В наборе JSONBench
    из публичного потока сети Bluesky (около 40~млн событий) при
    полной развёртке обнаруживается более 10~тысяч уникальных
    JSON-путей, при этом 99\% событий используют лишь около 50
    путей. Классическая СУБД с дисковыми страницами фиксированного
    размера не справляется с таблицей в 10~000 колонок.
    \item \textbf{Разрастание NULL-значений.} При $N~=~10^7$ строк и
    $5\cdot 10^3$ редких колонок только на битмаски валидности
    требуется порядка 6~ГБ.
    \item \textbf{Стоимость ALTER TABLE ADD COLUMN.} Каждый вызов
    добавления колонки требует $O(N)$ работы по инициализации
    значений по умолчанию для всей истории таблицы.
\end{itemize}

Теоретические основы колоночного хранения описаны в
C-Store~\cite{cstore}, MonetDB/X100~\cite{monetdb-x100},
Vectorwise~\cite{vectorwise}. Во всех этих работах рассматривается
случай фиксированной известной заранее схемы; случай динамической
схемы с произвольным ростом числа колонок не был предметом
систематического анализа до недавнего времени.

\section{Тип данных JSON в ClickHouse}

Наиболее продвинутым промышленным решением в области колоночного
хранения полуструктурированных данных на сегодня является новый
тип данных \texttt{JSON} в ClickHouse, введённый в версии~24.8 и
стабилизированный в~24.10~\cite{clickhouse-json}. Схема хранения
устроена следующим образом.

\textbf{Частично развёрнутая схема.} Тип декларируется как
\begin{verbatim}
JSON(
  max_dynamic_paths = N,
  field_1 Type_1,
  field_2 Type_2,
  ...
)
\end{verbatim}

Часть путей объявляется явно (с зафиксированным типом); они
становятся полноценными колонками основной таблицы со всеми
свойствами колоночного хранения (сжатие, zone~maps,
специализированные кодеки). Остальные пути помещаются в общее
хранилище \texttt{shared~data}.

\textbf{Shared~data.} Редкие пути хранятся внутри строки таблицы
как пара вложенных массивов: \texttt{\{"paths":[path\_1, \ldots],
"values":[value\_1, \ldots]\}}. По сути, это EAV-представление,
упакованное внутрь одной ячейки.

\textbf{Динамическое продвижение.} При вставке новых данных
система отслеживает частоту появления путей; пути, частота которых
превышает заданный порог \texttt{max\_dynamic\_paths} (по умолчанию
1024), автоматически <<продвигаются>> в полноценные колонки на
следующем слиянии.

Сильные стороны: автоматическое распознавание часто встречающихся
путей, эффективный колоночный доступ к популярным полям, компактное
хранение длинного хвоста.

Ограничения:
\begin{itemize}
    \item Ограниченный набор операций над \texttt{shared~data}:
    агрегации по путям оттуда существенно медленнее, чем по
    полноценным колонкам; индексы к ним не применяются.
    \item Отсутствие кросс-документной колоночной локальности для
    \texttt{shared~data}: чтобы узнать значение редкого поля в
    100~строках, ClickHouse должен разобрать 100~массивов.
    \item Стоимость слияний: перенос пути из \texttt{shared~data} в
    отдельную колонку требует переписывания всех кусков хранения
    (parts), затронутых полем.
\end{itemize}

Важно, что в текущей реализации ClickHouse редкие пути хранятся
внутри каждой строки, что означает дублирование ключей пути в
каждой строке (с применением словарной компрессии, но всё же в
логически миллионной кратности для часто встречающегося, но не
продвинутого пути).

\section{Выводы главы}

Проведённый анализ позволяет сформулировать требования к
разрабатываемому решению:
\begin{enumerate}
    \item Часто встречающиеся поля должны храниться как полноценные
    колонки~--- необходимое условие для достижения
    производительности, сопоставимой с ClickHouse, на типичных
    аналитических запросах.
    \item Редко встречающиеся поля должны храниться отдельно от
    основной таблицы: прямолинейное раздувание схемы неприемлемо;
    размещение редких значений внутри каждой строки сохраняет
    проблему чтения лишних данных при колоночных запросах.
    \item Должна существовать процедура автоматической миграции
    полей между режимами по мере роста их частоты, без
    необходимости ручного вмешательства в схему.
    \item Для разреженных полей нужно сохранить возможность
    применения стандартных операций (предикаты, индексы,
    агрегации)~--- желательно без специализированного кода, за
    счёт переиспользования существующих механизмов реляционного
    движка.
\end{enumerate}

Совокупно эти требования приводят к решению, описанному в
следующей главе: хранение каждого редкого пути в отдельной
side-таблице формата \texttt{(\_id, value)}, гибридное сканирование
одним специализированным оператором с подстановкой значений по
\texttt{\_id}, и явное разделение полей на закреплённые (pinned)
и разреженные (sparse) на этапе \texttt{CREATE TABLE}.

\finishrelatedwork

\chapterconclusion

Изложенные в настоящей главе подходы демонстрируют, что
существующие решения распределены между двумя крайностями:
документные системы теряют преимущества колоночного доступа,
классическое колоночное хранилище не справляется с произвольно
расширяющейся схемой. Гибридные подходы ClickHouse и Snowflake
частично закрывают этот разрыв, однако вносят собственные
ограничения. Предлагаемое в настоящей работе решение основано на
переиспользовании полноценных реляционных таблиц в качестве
хранилища для разреженных полей, что позволяет избежать
специализированного формата с ограниченной функциональностью.

\chapter{Проектирование решения}

\section{Общая идея}

Изучение существующих подходов привело к следующей комбинации
идей, которая и составляет архитектурное ядро настоящей работы.

\textbf{Основная таблица} (далее~--- main) играет роль источника
истины о множестве живых строк. Её обязательная колонка~---
служебный идентификатор \texttt{\_id BIGINT NOT NULL}, монотонно
выдаваемый при \textsc{insert}. В смешанной конфигурации main
дополнительно хранит \emph{закреплённые} (pinned) поля как обычные
физические колонки. Pinned-поля пользователь объявляет явно при
\texttt{CREATE TABLE} как поля, для которых заранее известна
высокая плотность заполнения; они получают полные преимущества
колоночного хранения (сжатие, zone~maps, векторизованные предикаты,
column~pruning).

\textbf{Вспомогательные таблицы} (далее~--- side) создаются по
одной на каждую пару (поле, тип), не объявленную pinned. Для каждой
такой пары заводится двухколоночная таблица вида:
\begin{verbatim}
CREATE TABLE db._dyn_<main>__<field>__<type> (
  _id BIGINT NOT NULL,
  value <type> NOT NULL
);
\end{verbatim}
В side-таблицу записываются только те строки, для которых исходное
поле было заполнено. Идентификатор \texttt{\_id} связывает строку
side-таблицы со строкой main, в которой это поле было встречено.

\textbf{Операция SELECT} реализуется специализированным
физическим оператором \texttt{scan\_computing\_table}, который
сканирует параллельно main и каждую запрошенную side-таблицу,
строит хеш-таблицу \texttt{\_id}~$\to$~позиция-в-выходе по
результатам main, и подставляет значения из side-таблиц в
соответствующие позиции; недостающие \texttt{\_id} остаются
NULL. Регулярные колонки, физически находящиеся в main,
читаются напрямую из её скана~--- без дополнительной операции
слияния. В терминах логической семантики это эквивалентно
LEFT~JOIN main и всех нужных side-таблиц по \texttt{\_id}, но без
создания цепочки JOIN-узлов в плане: операция гибридного сканирования
выполняется одним оператором.

\textbf{Эволюция схемы} происходит на стороне диспетчера: при
\textsc{insert} встретившаяся впервые пара (поле, тип) лениво
порождает соответствующую side-таблицу. Pinned-поля и их типы
зафиксированы при \texttt{CREATE TABLE}; динамическое продвижение
sparse-полей в main в настоящей работе не реализовано~--- секция
``Переход между режимами'' обсуждает, почему этот механизм отнесён
к направлениям будущей работы.

\section{Формализация модели}

Введём обозначения: $T$~--- main-таблица; $N$~--- число строк в
ней; $\mathcal{F}~=~\{f_1, \ldots, f_k\}$~--- множество
JSON-путей, встреченных хотя бы раз; $n_i$~--- число строк, в
которых поле $f_i$ было заполнено ненулевым значением.

Состояние каждого поля $f_i$ определяется одним из двух статусов:
\begin{itemize}
    \item \textbf{Pinned (закреплённое).} Указывается явно при
    \texttt{CREATE TABLE} через параметр \texttt{pinned\_columns}.
    Такое поле хранится в main-таблице как обычная физическая
    колонка, независимо от статистики заполненности. Pinned-поле
    одно для всех типов: при попытке вставить значение
    несовместимого с объявленным типом возвращается ошибка.
    \item \textbf{Sparse (разреженное).} По умолчанию все поля,
    не объявленные pinned, попадают сюда. Поле хранится только в
    отдельной side-таблице \texttt{\_dyn\_<main>\_\_<field>\_\_
    <type>}. Допускается несколько типов одного логического поля:
    каждая пара (поле, тип) получает свою side-таблицу.
\end{itemize}

В системе сохраняется один логический инвариант:
\begin{quote}\itshape
Если main-таблица содержит строку с идентификатором \texttt{\_id},
то на момент видимости этой строки все side-таблицы, описывающие
sparse-поля схемы, уже содержат соответствующие записи для этого
\texttt{\_id} (либо явное значение, либо его отсутствие).
\end{quote}
Из этого инварианта вытекает простое правило чтения: оператор
\texttt{scan\_computing\_table} использует main как ``якорь'' для
определения живых строк; видимая строка main гарантированно имеет
консистентное представление в side-таблицах. Side-таблицы могут
содержать ``висящие'' записи, не отражённые в main (например, из-за
обрыва транзакции между записью side и записью main),~--- такие
записи никогда не попадут в результат \textsc{select}, поскольку
\texttt{\_id} отсутствует в main и игнорируется при гибридном
сканировании. Подробное обсуждение этого инварианта и его роли
в обеспечении параллельности приведено в разделе
``Параллельность''.

\section{Обоснование формы хранения}

\subsection{Отдельная таблица вместо внутреннего формата}

Центральное архитектурное решение~--- хранить разреженные поля в
полноценных двухколоночных таблицах, а не в специальном внутреннем
формате (как \texttt{shared~data} в ClickHouse). Аргументы в пользу
этого выбора:

\begin{itemize}
    \item \textbf{Переиспользование всей реляционной машинерии.}
    Любая оптимизация, применяемая к основной таблице
    (zone~maps, векторизация, projection~pushdown, индексы,
    сжатие), автоматически применима и к вспомогательной таблице.
    Не требуется писать специализированный код для разреженного
    хранилища.
    \item \textbf{Совместимость с MVCC и WAL.} Операции над
    вспомогательной таблицей проходят по тому же пути, что и над
    основной, автоматически получая транзакционные гарантии и
    журналирование.
    \item \textbf{Простота отладки и диагностики.}
    Вспомогательная таблица~--- это обычная таблица; её можно
    посмотреть отдельным \textsc{select}, проверить содержимое,
    оценить размер, применить EXPLAIN.
    \item \textbf{Возможность независимой оптимизации.}
    Вспомогательная таблица может иметь собственный индекс по
    \texttt{\_id} (например, zone~map по \texttt{\_id} полезен
    при выборочных предикатах по основной таблице с последующим
    JOIN).
\end{itemize}

\subsection{Логический идентификатор вместо физической позиции строки}

Альтернатива~--- использовать физический номер строки
(row~position) как ключ связи. Отбрасывается по причинам
устойчивости к операциям над основной таблицей (физическая позиция
может меняться при сжатии, ребалансировке row~groups, перестройке
индексов; логический \texttt{\_id} остаётся стабильным) и
отсутствия кросс-зависимостей (вспомогательная таблица не должна
знать о физическом расположении строк в основной; это упрощает
параллельную обработку~--- \textsc{insert} в основную и
вспомогательные таблицы могут идти одновременно, как только известен
\texttt{\_id}).

Идентификаторы \texttt{\_id} генерируются монотонно: для каждой
операции \textsc{insert} система запрашивает у основной таблицы
текущее число строк, присваивает вставляемым строкам
последовательный диапазон $[\texttt{start\_id},
\texttt{start\_id}+|C|)$, и этот диапазон далее используется как
ключ связи.

\subsection{Специализированный оператор гибридного сканирования}

Чтение данных с разреженными полями реализовано отдельным физическим
оператором \texttt{scan\_computing\_table}, а не цепочкой LEFT~JOIN
в логическом плане. Семантически оба варианта эквивалентны, но
выделение в отдельный оператор имеет ряд преимуществ:

\begin{itemize}
    \item \textbf{Локальность алгоритма.} Параллельное чтение
    main и всех нужных side-таблиц, построение хеш-таблицы
    \texttt{\_id}~$\to$~позиция и подстановка значений происходят
    внутри одного оператора. Это позволяет не материализовывать
    промежуточные представления отдельных таблиц.
    \item \textbf{Естественная интеграция с column~pruning.}
    Оператор принимает список запрошенных колонок и обращается
    только к тем side-таблицам, поля которых упомянуты в запросе.
    Пары (поле, тип), не упомянутые в запросе, не сканируются
    вовсе.
    \item \textbf{Единый источник истины.} Множество <<живых>>
    \texttt{\_id} определяется одним сканом main; для каждой
    side-таблицы строится отдельный проход с подстановкой
    значений по уже построенной хеш-таблице. Не возникает
    проблем согласованности нескольких независимых
    JOIN-операторов.
    \item \textbf{Совместимость с параллельностью.} Из инварианта
    раздела ``Формализация'' следует, что сканирование main и
    side-таблиц можно вести параллельно, не дожидаясь окончания
    параллельных \textsc{insert}: видна будет только та часть
    данных, которая зафиксирована в main на момент начала
    транзакции (см. раздел ``Параллельность'').
\end{itemize}

Альтернатива в виде LEFT~JOIN на уровне логического плана была
отброшена: она вынуждает оптимизатор и валидатор учитывать
служебные узлы JOIN, что увеличивает поверхность изменений в
зрелых компонентах кодовой базы, не давая дополнительной
функциональности~--- алгоритмически операция остаётся такой же.

\section{Когда стоит делать колонку pinned}

В предложенной модели решение о размещении поля принимается
однократно на этапе \texttt{CREATE TABLE} и далее не меняется.
Аналитик объявляет некоторый список pinned-полей (возможно, пустой);
все остальные поля автоматически становятся sparse. Этот выбор
влияет на три аспекта производительности:

\begin{itemize}
    \item \textbf{Скорость \textsc{select} по полю.} Чтение
    pinned-поля происходит напрямую из колоночного буфера
    main-таблицы и пользуется всей колоночной машинерией
    (zone~maps, векторизация, сжатие). Чтение sparse-поля
    требует прохода по соответствующей side-таблице с
    хеш-подстановкой по \texttt{\_id}~--- это дороже на
    константный множитель, но всё ещё линейно по числу строк.
    \item \textbf{Стоимость хранения редких полей.} Pinned-поле
    выделяет в main-таблице буфер на все $N$ строк независимо от
    плотности заполнения. Для редкого поля это пустая трата
    памяти. Sparse-поле занимает место только под фактически
    заполненные строки.
    \item \textbf{Стоимость эволюции схемы.} Pinned-поля
    фиксированы при \texttt{CREATE TABLE}; добавление нового
    pinned-поля задним числом не предусмотрено. Sparse-поля
    добавляются лениво в момент первой вставки без блокировок и
    миграции.
\end{itemize}

Эмпирическое правило: если поле упоминается в большинстве запросов
и заполнено в более чем 30--50\% строк~--- кандидат на pinned. Поля
с плотностью ниже 10\% разумно оставлять sparse независимо от
частоты использования.

Альтернативный механизм автоматического продвижения sparse-поля в
main по достижении порога заполненности рассматривался на ранней
стадии работы и отнесён к будущим направлениям; обсуждение
см.~раздел ``Переход между режимами''.

\section{Переход между режимами}

Автоматическое продвижение sparse-поля в pinned-колонку
main-таблицы по достижении порога заполненности рассматривалось на
ранней стадии работы. В предложенной архитектуре от него решено
отказаться по совокупности причин.

\textbf{Сложность под нагрузкой.} Продвижение~--- $O(N)$ операция
(добавление физической колонки в main с инициализацией значений по
умолчанию + полное сканирование side-таблицы для миграции
значений). Если продвижение запускается в фоне параллельно с
\textsc{insert}, требуется аккуратная синхронизация: новые
\textsc{insert} могут писать в side, ещё не дочитанную миграцией.
Если выполняется блокирующим образом, продвижение приостанавливает
все \textsc{insert} в данное поле на время операции.

\textbf{Нарушение инварианта во время миграции.} Главное достоинство
текущей схемы~--- стабильное правило ``main определяет видимость, side
дополняет значения''. Промежуточное состояние ``часть строк уже в
main-колонке, часть ещё в side-таблице'' нарушает этот инвариант и
требует двойного источника правды на время миграции, что
существенно усложняет реализацию \texttt{scan\_computing\_table}.

\textbf{Необходимость обратного перехода.} Если продвижение
осуществляется автоматически по статистике, симметрично следует
поддерживать и понижение (удаление физической колонки при падении
плотности). Эта операция $O(N)$ и блокирующая по той же причине,
что и прямой переход.

В предложенной архитектуре выбран альтернативный путь: статус поля
(pinned либо sparse) фиксируется при \texttt{CREATE TABLE} и в
дальнейшем не меняется. Это позволяет сохранить инвариант
``main~--- источник истины'' и упростить и оператор сканирования, и
анализ параллельной семантики (см.~раздел ``Параллельность'').
Автоматическое продвижение рассматривается как направление будущей
работы в случаях, когда статистика заполненности неизвестна на
этапе \texttt{CREATE TABLE} и заранее зафиксировать список pinned
невозможно.

\section{Плюсы и минусы подхода}

\textbf{Преимущества:}
\begin{enumerate}
    \item Колоночная производительность на pinned-полях: они
    физически лежат в main и пользуются всей колоночной
    машинерией (zone~maps, векторизация, сжатие, column~pruning).
    \item Компактное хранение редких полей: side-таблицы содержат
    только фактически заполненные строки; неупомянутые в запросе
    side-таблицы не сканируются вовсе.
    \item Автоматическая эволюция схемы: не требуется ручного
    переопределения структуры при появлении новых sparse-полей~---
    side-таблица для новой пары (поле, тип) создаётся лениво при
    первой вставке.
    \item Отсутствие специальных структур данных: side-таблица~---
    это обычная колоночная таблица; вся транзакционная и
    инфраструктурная машинерия (MVCC, WAL, индексы) применима
    без модификаций.
    \item Стабильная семантика \texttt{\_id}: внешние индексы и
    ссылки на строки устойчивы к внутренним операциям движка
    хранения.
    \item Простой инвариант параллельности (раздел
    ``Параллельность''): если строка видна в main, её
    представление в side-таблицах гарантированно консистентно.
\end{enumerate}

\textbf{Недостатки и компромиссы:}
\begin{enumerate}
    \item Чтение sparse-поля дороже чтения pinned-колонки:
    хеш-подстановка по \texttt{\_id} вносит накладные расходы
    порядка $O(n_i)$ операций хеширования на запрос.
    \item Отсутствие автоматического продвижения: при плохом
    выборе списка pinned-полей часто запрашиваемые
    sparse-поля сохраняют повышенную стоимость чтения. Решение
    отнесено к направлениям будущей работы.
    \item Число side-таблиц может расти значительно: при
    10\,000 уникальных путей в каталоге возникает 10\,000
    физических таблиц; каждая потребляет некоторый объём
    памяти на метаданные.
\end{enumerate}

\section{Псевдокод INSERT}

Для формализации опишем \textsc{insert} в таблицу с вычисляемой
схемой в виде псевдокода (листинг~\ref{lst:insert}). Алгоритм
выполняется одним физическим оператором
\texttt{operator\_insert\_computing} в исполнителе под одной
транзакцией, что гарантирует атомарность по отношению к
параллельным запросам.

\begin{algorithm}[!h]
\caption{Алгоритм \textsc{insert} в таблицу с вычисляемой схемой}
\label{lst:insert}
\begin{algorithmic}
  \Function{Insert}{row-chunk $C$, схема $S$, main-таблица $T$}
    \State \texttt{// pre-flight DDL на стороне диспетчера}
    \ForAll{колонок $c_j$ из $C$}
      \State $S.\mathrm{append}(f_j, t_j)$
      \If{пара $(f_j, t_j)$ ещё не имеет side-таблицы и не
        pinned}
        \State создать side-таблицу
          \texttt{\_dyn\_<T>\_\_<$f_j$>\_\_<$t_j$>}
      \EndIf
    \EndFor
    \State \texttt{// выполнение в исполнителе под одной транзакцией}
    \State $\mathtt{start\_id} \gets$ число строк $T$
    \State $\mathtt{main\_chunk} \gets [\_id]$ + значения
      pinned-колонок из $C$
    \ForAll{колонок $c_j$ из $C$, не являющихся pinned}
      \State построить $\mathtt{side\_chunk}_j$ из пар
        $(\mathtt{start\_id}+i, c_j[i])$ по ненулевым строкам
    \EndFor
    \State \texttt{// первая фаза: side-таблицы}
    \ForAll{$\mathtt{side\_chunk}_j$}
      \State \texttt{storage\_append} в side-таблицу
        $(f_j, t_j)$ с $\mathit{txn\_id}$
    \EndFor
    \State \texttt{// вторая фаза: main-таблица}
    \State \texttt{storage\_append} в $T$ с
      $\mathtt{main\_chunk}$ и $\mathit{txn\_id}$
    \State \texttt{// фиксация}
    \State $\mathit{commit\_id} \gets
      \mathrm{txn\_manager}.\textsc{commit}()$
    \State \texttt{storage\_commit\_append} для каждой
      затронутой таблицы (main и side) с $\mathit{commit\_id}$
  \EndFunction
\end{algorithmic}
\end{algorithm}

\textbf{Принципиальный порядок записи: side первыми, main
последней.} Это поддерживает инвариант ``main определяет
видимость, side хранит данные''. Если транзакция завершилась
успешно, все side-записи с $\mathit{commit\_id}$ становятся видимыми
одновременно с main-записью с тем же $\mathit{commit\_id}$~--- их
снимки атомарны. Если транзакция оборвалась после записи в
side-таблицу, но до записи в main, side-записи остаются помечены
исходным $\mathit{txn\_id}$ (без $\mathit{commit\_id}$) и невидимы
любому читателю. Поскольку \texttt{\_id} такой строки в main также
отсутствует, оператор \texttt{scan\_computing\_table} никогда не
обратится к этой строке, даже если бы видел её.

Алгоритмическая модель \textsc{select} симметрична: один
\texttt{storage\_scan} по main даёт множество живых \texttt{\_id} и
значения pinned-колонок, дальше каждая упомянутая в запросе
sparse-side сканируется один раз с подстановкой значений в
позиции, найденные по хеш-таблице
\texttt{\_id}~$\to$~позиция-в-выходе.

\section{Параллельность}\label{sec:parallel}

Корректность параллельного выполнения операций над таблицей с
вычисляемой схемой обеспечивается не на уровне внешней сериализации
(хеш-маршрутизацией ``одна коллекция $\to$ один исполнитель'',
которая на текущем этапе используется в реализации, но в будущем
будет ослаблена для повышения масштабируемости), а на уровне самого
устройства схемы: \emph{инвариант ``main~--- источник истины''}
делает параллельное чтение и запись безопасными независимо от
порядка планирования транзакций.

\subsection{Базовый инвариант и его следствия}

Напомним инвариант, введённый в разделе ``Формализация модели'':
\begin{quote}\itshape
Если строка с идентификатором \texttt{\_id} видна читателю в main,
то соответствующие записи во всех side-таблицах либо уже видны тому
же читателю, либо явно отсутствуют (что трактуется как NULL).
\end{quote}

Инвариант поддерживается на этапе записи через порядок
\emph{side $\to$ main}: исполнитель сначала проводит
\texttt{storage\_append} во все side-таблицы под текущим
$\mathit{txn\_id}$, затем~--- в main, и только потом получает у
менеджера транзакций общий $\mathit{commit\_id}$ и применяет
\texttt{storage\_commit\_append} ко всем затронутым таблицам с этим
$\mathit{commit\_id}$. MVCC-протокол хранилища гарантирует, что
никакой читатель не видит записи с $\mathit{txn\_id}$ (только
$\mathit{commit\_id}$, выданный после успешного коммита), поэтому
запись становится видимой во всех таблицах одновременно.

При оборванной транзакции (отказ процесса, явный
\texttt{rollback}) ни одна из записей не получает
$\mathit{commit\_id}$~--- они остаются невидимыми. ``Висящие''
записи в side-таблицах безопасно игнорируются последующим
\texttt{scan\_computing\_table}, поскольку без соответствующей
живой строки в main их \texttt{\_id} не попадает в множество ключей
для хеш-подстановки.

\subsection{Параллельный INSERT}

Рассмотрим две параллельные транзакции $\tau_A$ и $\tau_B$, каждая
выполняющая \textsc{insert} нескольких строк в одну и ту же
вычисляемую таблицу.

Каждая транзакция вначале запрашивает у менеджера хранилища текущее
число строк main как $\mathtt{start\_id}$. Эта операция атомарна на
уровне акторного менеджера хранилища: его mailbox обрабатывает
сообщения последовательно, поэтому два параллельных запроса
получают разные значения, отражающие фактический порядок их
обработки. Соответственно $\tau_A$ и $\tau_B$ получают
непересекающиеся диапазоны \texttt{\_id}.

Дальше каждая транзакция действует независимо. Записи $\tau_A$ в
side-таблицы помечены $\mathit{txn\_id}_A$ и невидимы $\tau_B$ (и
наоборот). Когда $\tau_A$ применяет \texttt{storage\_commit\_append}
со своим $\mathit{commit\_id}_A$, все её записи становятся видимыми
последующим транзакциям с большим $\mathit{start\_time}$. Записи
$\tau_B$, ещё не закоммиченные, остаются невидимыми. Никакого
взаимодействия между $\tau_A$ и $\tau_B$ нет: они работают с
непересекающимися \texttt{\_id} и не делят ни одной строки.

\subsection{Параллельный SELECT и INSERT}

Параллельный \textsc{select} в момент незавершённой $\tau_A$
открывает свою транзакцию $\tau_S$ с собственным
$\mathit{start\_time}$. Сканирование main отдаёт только те строки,
$\mathit{commit\_id}$ которых не превышает $\mathit{start\_time}$
читателя~--- стандартный snapshot isolation. Если $\tau_A$ ещё не
закоммитилась к моменту начала $\tau_S$, её записи не видны ни в
main, ни в side. Если $\tau_A$ закоммитилась раньше начала
$\tau_S$, её записи видны и там, и там одновременно~--- инвариант
``main определяет видимость''.

Случай, когда $\tau_A$ закоммитилась в момент исполнения $\tau_S$
(между сканом main и сканом side), требует уточнения. MVCC-снимок
фиксируется при \texttt{begin\_transaction}, не при каждом
обращении к хранилищу,~--- поэтому даже если физически сканы main и
side происходят в разное время, оба видят одно и то же состояние,
консистентное с $\mathit{start\_time}$ читателя. ``Подоспевшие''
изменения $\tau_A$ не попадают в результат $\tau_S$, что и требуется
от snapshot~isolation.

\subsection{Параллельный DELETE}

\textsc{delete} строк с заданным WHERE реализуется отдельным
оператором \texttt{operator\_delete\_computing} и выполняется в
строго обратном порядке: \emph{main первым, side после}.

Сначала исполнитель помечает в main все строки, попадающие под
WHERE, как удалённые (под $\mathit{txn\_id}$); затем~--- те же
строки в каждой side-таблице. Когда транзакция коммитится, помеченные
``удалённые'' записи получают $\mathit{commit\_id}$ и перестают быть
видимы новым читателям. Симметрия с \textsc{insert}: видимость
``удалённого состояния'' возникает одновременно во всех таблицах.

Параллельный \textsc{select} в момент незавершённого DELETE по
snapshot~isolation видит либо все строки целиком (если его
$\mathit{start\_time}$ предшествует $\mathit{commit\_id}$ удаления),
либо ни одной из удаляемых (после коммита). Никаких ``половинных''
состояний~--- то есть видимая в main строка без значений из side
или, наоборот, висящие в side значения без живой строки в main~---
никогда не наблюдается.

Параллельный \textsc{insert} в момент незавершённого DELETE
работает с непересекающимися \texttt{\_id}, поскольку
\texttt{start\_id} для нового \textsc{insert} больше любого
существующего \texttt{\_id}, а удаляемые \texttt{\_id}
заведомо меньше.

\subsection{Параллельный UPDATE}

\textsc{update} реализован как ``re-alloc с полным копированием'':
оператор \texttt{operator\_update\_computing} вычитывает строки,
попадающие под WHERE, применяет SET-выражения к их значениям,
выделяет новый диапазон \texttt{\_id} через
\texttt{storage\_total\_rows}, делает \textsc{insert} в side и main
с новыми \texttt{\_id} и \textsc{delete} старых строк в main и
side. Всё это в одной транзакции.

Параллельный \textsc{select} в момент незавершённого UPDATE по
снимку видит либо старую версию строки (если до коммита), либо
новую (если после). Никогда не видит обе одновременно, поскольку и
вставка ``нового'' и удаление ``старого'' получают один
$\mathit{commit\_id}$.

Параллельные UPDATE/INSERT/DELETE по \emph{разным} \texttt{\_id}
работают полностью независимо: операции над диапазонами с разными
\texttt{\_id} не делят ни одной строки и потому не имеют конфликтов
данных.

Параллельные UPDATE по \emph{одному и тому же} логическому
\texttt{\_id} (с пересекающимися WHERE-условиями) образуют
write-write конфликт. Otterbrix~--- как и большинство классических
аналитических колоночных систем~--- не обнаруживает такие конфликты
автоматически; разрешение в таких случаях соответствует
last-writer-wins по $\mathit{commit\_id}$. Это ограничение
существующей MVCC-реализации, не связанное со схемой вычисляемой
таблицы; устранение требует доработки на уровне всего хранилища.

\subsection{Безопасность по отношению к будущему ослаблению
сериализации}

Текущая реализация хеш-маршрутизирует запросы к одной коллекции на
один акторный исполнитель, что само по себе сериализует все
операции с этой коллекцией. Однако такое ограничение является
артефактом текущего архитектурного решения, не свойством схемы
вычисляемой таблицы. В будущем планируется балансировать запросы по
исполнителям независимо от коллекции (например, по сессии или
по размеру очереди), что увеличит масштабируемость, но потребует
полагаться только на свойства самой схемы для корректности.

Все приведённые выше рассуждения о параллельности \emph{не зависят}
от хеш-маршрутизации: они полагаются исключительно на
MVCC-механизмы хранилища и порядок записи ``side $\to$ main''
(симметрично для удаления ``main $\to$ side''). Таким образом,
предложенная архитектура естественно совместима с будущим
ослаблением сериализации без дополнительных модификаций.

\chapterconclusion

Предложенная модель хранения сочетает сильные стороны колоночного
представления (для часто используемых pinned-полей) и разреженного
EAV (для редких полей), обходя ключевые ограничения каждого из
подходов. Центральные архитектурные решения~--- хранение
sparse-полей в отдельных двухколоночных таблицах, связанных через
стабильный логический идентификатор \texttt{\_id}, гибридное
сканирование одним специализированным оператором, и инвариант
``main~--- источник истины'' для согласованности параллельных
операций~--- позволяют достичь приемлемого баланса между стоимостью
хранения, стоимостью чтения и сложностью реализации, а также
естественной совместимости с будущим ослаблением сериализации на
уровне исполнителей.

\chapter{Реализация}\label{ch:impl}

Настоящая глава описывает реализацию предложенной в главе~3 схемы
хранения в кодовой базе Otterbrix. Основная цель раздела~--- показать,
как архитектурные решения отображаются на конкретные компоненты
кода, и какие структурные изменения потребовалось внести в
различные подсистемы.

\section{Обзор затронутых компонентов}

Реализация задействует следующие ключевые подсистемы Otterbrix:

\begin{itemize}
    \item \texttt{components/catalog/}~--- метаданные таблиц.
    Введён новый класс \texttt{computed\_schema}, отвечающий за
    учёт зарегистрированных полей и их типов.
    \item \texttt{components/physical\_plan/operators/}~---
    физические операторы. Введены четыре новых оператора:
    \texttt{scan\_computing\_table} (гибридное сканирование),
    \texttt{operator\_insert\_computing},
    \texttt{operator\_delete\_computing},
    \texttt{operator\_update\_computing} (multi-collection DML под
    одной транзакцией).
    \item \texttt{services/collection/}~--- исполнитель планов.
    Цикл DML расширен поддержкой ``вычисляемых''
    DML-операторов: \texttt{intercept\_dml\_io\_} получил отдельные
    ветви для \texttt{insert\_computing},
    \texttt{remove\_computing} и \texttt{update\_computing},
    которые ведут многотабличный жизненный цикл транзакции (запись
    в N+1 таблиц, WAL physical, общий commit с одним
    $\mathit{commit\_id}$, общий COMMIT marker).
    \item \texttt{services/dispatcher/}~--- центральная точка
    маршрутизации запросов. Pre-flight для \textsc{insert} в
    вычисляемую таблицу создаёт недостающие side-таблицы и
    проставляет на узел плана список нужных side для последующего
    исполнителя.
    \item \texttt{services/disk/}~--- менеджер дискового
    хранилища. Добавлена поддержка создания, удаления и сканирования
    side-таблиц (как обычных регулярных таблиц).
    \item \texttt{components/sql/transformer/}~--- SQL-трансформер.
    Добавлена поддержка синтаксиса \texttt{CREATE TABLE db.t~();}
    (вычисляемая схема), \texttt{WITH (pinned\_columns~=~\ldots)}
    (список pinned-полей) и \texttt{INSERT \ldots\ VALUES (json('\ldots'))}.
    \item \texttt{components/logical\_plan/}~--- узлы логического
    плана. На \texttt{node\_create\_collection\_t} добавлен флаг
    динамической схемы и опциональный список pinned-полей; на
    \texttt{node\_insert\_t}, \texttt{node\_delete\_t} и
    \texttt{node\_update\_t}~--- метаданные \texttt{computing\_sides}
    для маршрутизации в специализированные операторы.
    \item \texttt{components/table/}~--- низкоуровневое колоночное
    хранилище. Без существенных модификаций~--- side-таблицы
    хранятся как обычные регулярные таблицы.
    \item \texttt{components/vector/}~--- операции над векторами
    значений. Добавлены placeholder-векторы (без выделения буфера
    для непроецированных колонок) для поддержки column~pruning в
    \texttt{scan\_computing\_table}.
\end{itemize}

\section{Класс представления вычисляемой схемы}

Класс \texttt{computed\_schema}
(\texttt{components/catalog/computed\_schema.hpp}/.cpp)
инкапсулирует учёт метаданных таблицы с вычисляемой схемой.

\subsection{Поля состояния}

Класс \texttt{computed\_schema} хранит следующее состояние:

\begin{itemize}
    \item отображение имени поля в перечень типов, в которых оно
    встречалось;
    \item упорядоченный список пар <<имя поля~--- тип>> в порядке
    их первого появления;
    \item множество имён полей, объявленных pinned при
    \texttt{CREATE TABLE} (если таковые есть);
    \item флаг \texttt{is\_sparse}, указывающий, является ли
    таблица таблицей с вычисляемой (разреженной) схемой; для обычных
    регулярных таблиц этот флаг равен \texttt{false} и логика
    side-таблиц отключена.
\end{itemize}

Динамическое отслеживание статистики заполненности и счётчики
non-null убраны вместе с механизмом автоматического продвижения
(см.~раздел ``Переход между режимами'' в главе~3).

Центральное понятие~--- имя side-таблицы, формируемое статической
функцией \texttt{side\_table\_name(main, field, type)} по правилу
\texttt{\_dyn\_<main>\_\_<field>\_\_<type\_enum>}. Префикс
\texttt{\_dyn\_} зарезервирован для служебных таблиц и гарантирует,
что автоматически порождённые имена не конфликтуют с
пользовательскими. Тип в суффиксе необходим потому, что одно и то
же логическое поле может существовать в нескольких типах: если
первый \textsc{insert} вставил \texttt{\{"price":~100\}}, а
следующий~--- \texttt{\{"price":~"free"\}}, в схеме появляются две
отдельные side-таблицы \texttt{\_dyn\_t\_\_price\_\_INT64} и
\texttt{\_dyn\_t\_\_price\_\_STRING}.

\subsection{Центральные операции}

\begin{description}
    \item[\texttt{append(field\_name, type)}] Регистрирует новую
    пару (field, type) в схеме, если её ещё нет. Идемпотентна.
    \item[\texttt{latest\_types\_struct()}] Возвращает текущую
    виртуальную схему таблицы в виде структурного типа: каждая
    зарегистрированная пара (field, type) фигурирует как поле.
    Используется валидатором для проверки запросов на
    soundness относительно текущей схемы.
    \item[\texttt{side\_table\_name(main, field, type)}]
    Генерирует имя side-таблицы. Статическая функция; не
    обращается к состоянию объекта.
\end{description}

\subsection{Потокобезопасность}

\texttt{computed\_schema} синхронизируется внешним блокирующим
мьютексом на уровне каталога. Все изменения схемы
(\texttt{append}) выполняются под мьютексом каталога во время
pre-flight \textsc{insert}. Поскольку изменения~--- лишь добавление
новых пар (field, type), а не модификация существующих, конкуренция
на мьютексе минимальна. Чтения схемы (\texttt{latest\_types\_struct})
требуются только при валидации запросов и физически не конфликтуют
ни с какими записями.

\section{Диспетчер: pre-flight для DML}

Обработчик каждого DML-узла (\textsc{insert}, \textsc{delete},
\textsc{update}) в диспетчере, обнаружив таблицу с вычисляемой
схемой (\texttt{catalog\_.table\_computes(id)}), выполняет короткую
pre-flight-фазу под мьютексом диспетчера и далее передаёт работу в
исполнитель:

\begin{itemize}
    \item Для \textsc{insert}: проход по колонкам входного чанка с
    регистрацией пары (поле, тип) в \texttt{computed\_schema} и
    созданием соответствующих side-таблиц для тех пар, что
    впервые увидены (атомарно через
    \texttt{manager\_disk\_t::create\_storage\_with\_columns},
    \texttt{register\_collection},
    \texttt{catalog\_.create\_table}). Список созданных и уже
    существующих side-таблиц проставляется на узел
    \texttt{node\_insert\_t} в виде поля
    \texttt{computing\_sides}.
    \item Для \textsc{delete} и \textsc{update}: текущая схема
    таблицы (полный список side) проставляется на соответствующий
    узел~--- без DDL-операций, поскольку DML только модифицирует
    существующие записи.
\end{itemize}

После pre-flight план отправляется в исполнитель стандартным
вызовом \texttt{execute\_plan\_impl}~--- ровно тот же путь, что и
для обычной DML на регулярной таблице. Мьютекс диспетчера
освобождается сразу после успешной маршрутизации, и продолжительная
DML-работа над main и side-таблицами выполняется на исполнителе
без удержания глобального мьютекса.

\section{Исполнитель: операторы вычисляемой DML}

Главное архитектурное изменение реализации~--- перенос всей
многотабличной DML-работы из диспетчера в специализированные
физические операторы в исполнителе. Это позволяет проводить запись
в main и все нужные side-таблицы под одной транзакцией с одним
$\mathit{txn\_id}$ и одним $\mathit{commit\_id}$, что даёт
real-MVCC-атомарность.

\subsection{operator\_insert\_computing}

Минималистичный read\_write-оператор: хранит имя main-таблицы и
список side (\texttt{computing\_sides}), полученный с
\texttt{node\_insert\_t}. В \texttt{on\_execute\_impl} переносит
вывод левого ребёнка (исходный \texttt{data\_chunk\_t} от
\texttt{operator\_data}) в собственный \texttt{output\_} и
переходит в режим \texttt{async\_wait}. Вся работа происходит в
ветке \texttt{intercept\_dml\_io\_::case insert\_computing} в
исполнителе:

\begin{enumerate}
    \item \texttt{storage\_total\_rows} на main~--- получение
    диапазона новых \texttt{\_id}.
    \item Разбиение пользовательского чанка: pinned-колонки и
    \texttt{\_id} формируют чанк для main, каждое sparse-поле
    формирует двухколоночный (\texttt{\_id}, value) чанк для своей
    side-таблицы (только ненулевые строки).
    \item \texttt{storage\_append} на каждую side-таблицу с
    текущим $\mathit{txn\_id}$; затем \texttt{storage\_append} на
    main с тем же $\mathit{txn\_id}$.
    \item WAL physical записи для каждой затронутой таблицы.
    \item \texttt{txn\_manager.commit} $\to$
    $\mathit{commit\_id}$;
    \texttt{storage\_commit\_append} для каждой таблицы с этим
    $\mathit{commit\_id}$.
    \item Один WAL COMMIT marker для всей транзакции.
\end{enumerate}

При ошибке на любом шаге выполняется
\texttt{storage\_revert\_append} для уже записанных таблиц и
\texttt{txn\_manager.abort}. Поскольку $\mathit{commit\_id}$ не
был выставлен, эти записи никому невидимы. Side-таблицы могут
содержать ``висящие'' уже-помеченные-удалёнными записи; они
безопасно отчищаются последующим garbage collection.

\subsection{operator\_delete\_computing}

В отличие от обычного \texttt{operator\_delete}, который работает
с одной таблицей, computing-вариант знает обо всех side-таблицах
схемы. Левое поддерево оператора~---
\texttt{scan\_computing\_table~+~operator\_match}: оно отдаёт
чанки с логическими \texttt{\_id} в \texttt{chunk.row\_ids} для
тех строк, что удовлетворяют WHERE.

В ветке \texttt{intercept\_dml\_io\_::case remove\_computing}:

\begin{enumerate}
    \item Собирается множество удаляемых \texttt{\_id}.
    \item Для каждой целевой таблицы (main и каждая side) делается
    \texttt{storage\_scan\_batched} с подбором физических позиций
    тех строк, чьё значение колонки \texttt{\_id} принадлежит
    собранному множеству. Это эквивалентно поиску по индексу
    \texttt{\_id}, если такой индекс есть, или линейному скану с
    pruning по zone~maps в противном случае.
    \item \texttt{storage\_delete\_rows} с найденными физическими
    позициями и текущим $\mathit{txn\_id}$ на каждую таблицу. Main
    помечается \emph{первой}, side~--- после, согласно инварианту
    ``main определяет видимость, удаление сначала отменяет видимость
    в main, потом чистит side''.
    \item WAL physical\_delete, \texttt{txn\_manager.commit},
    \texttt{storage\_commit\_delete} на каждую таблицу с общим
    $\mathit{commit\_id}$, и единственный WAL COMMIT marker.
\end{enumerate}

\subsection{operator\_update\_computing}

Реализует ``re-alloc с полным копированием'': для каждой
матчащейся строки выделяется новый \texttt{\_id}, под этот новый
\texttt{\_id} записываются все поля строки (с применением
SET-выражений), и старая строка удаляется. Структурно объединяет
ветви \texttt{insert\_computing} и \texttt{remove\_computing}.

Левое поддерево, аналогично DELETE, состоит из
\texttt{scan\_computing\_table~+~operator\_match}, но проектирует
\emph{все} колонки виртуальной схемы (не только row\_ids), чтобы
получить полные значения матчащихся строк. SET-выражения
применяются in-place к выходному чанку через
\texttt{update\_expr\_t::execute}, после чего:

\begin{enumerate}
    \item \texttt{storage\_total\_rows} $\to$ новый диапазон
    \texttt{\_id}.
    \item Сборка новых per-side и main-чанков, аналогично
    \texttt{insert\_computing}.
    \item \texttt{storage\_append} в side, затем в main.
    \item Сбор физических позиций старых \texttt{\_id} в каждой
    таблице (сканированием с фильтрацией).
    \item \texttt{storage\_delete\_rows} на тех же таблицах под
    тем же $\mathit{txn\_id}$.
    \item WAL physical для всех вставок и удалений,
    \texttt{txn\_manager.commit},
    \texttt{storage\_commit\_append}+\texttt{storage\_commit\_delete}
    на каждую таблицу с одним $\mathit{commit\_id}$, общий WAL
    COMMIT marker.
\end{enumerate}

Под одним $\mathit{commit\_id}$ становятся видны одновременно и
``новая'' строка, и снятие видимости со ``старой''~--- результат
наблюдаем как атомарная подмена.

\section{Оператор scan\_computing\_table}

Гибридное сканирование для \textsc{select} реализовано отдельным
физическим оператором
\texttt{components/physical\_plan/operators/scan/scan\_computing\_table.\{hpp,cpp\}}.

Оператор принимает на вход:
\begin{itemize}
    \item имя main-таблицы;
    \item список side (имя~--- тип значения), параллельный позиции
    в виртуальной схеме (без main);
    \item номер pinned-колонок main (доступен через
    \texttt{computed\_schema});
    \item список запрошенных позиций виртуальной схемы
    (\texttt{projected\_cols}~--- стандартный механизм
    column~pruning).
\end{itemize}

Алгоритм \texttt{await\_async\_and\_resume}:

\begin{enumerate}
    \item \emph{Параллельный запуск сканов.} Отправляются
    одновременно асинхронные сообщения
    \texttt{storage\_scan\_batched} к менеджеру диска для main и
    для каждой запрошенной side-таблицы. \texttt{co\_await}
    собирает все ответы; благодаря акторной модели Otterbrix
    физическое чтение разных таблиц может идти на разных
    исполнителях диска параллельно.
    \item \emph{Сканирование main.} Выходные чанки main отдают
    значения pinned-колонок напрямую (они физически здесь же), а
    также набор живых \texttt{\_id} (записываются в
    \texttt{chunk.row\_ids} выходных чанков как канал
    мета-идентификаторов, доступный последующим операторам:
    operator\_match, operator\_delete\_computing,
    operator\_update\_computing).
    \item \emph{Построение хеш-таблицы
    \texttt{\_id}~$\to$~позиция-в-выходе.} Один проход по
    собранным \texttt{\_id} строит словарь.
    \item \emph{Сканирование side с подстановкой.} Для каждой
    запрошенной sparse-side: проход по её чанкам, для каждой
    строки $(\_id, value)$ ищется позиция-в-выходе в хеш-таблице,
    \texttt{value} записывается в соответствующую колонку
    выходного чанка по этой позиции. Строки без записи в side
    остаются NULL (выходные векторы инициализированы как
    ``все~NULL'').
\end{enumerate}

Сложность по времени: $O(N + \sum_i n_i)$, где $N$~--- число
живых строк main, $n_i$~--- число строк в $i$-й запрошенной
side-таблице. Сложность по памяти: $O(N)$ для хеш-таблицы плюс
$O(N \cdot k)$ для выходных буферов, где $k$~--- число запрошенных
колонок (стандартно). Side-таблицы, не упомянутые в запросе, не
сканируются вовсе.

\section{SQL-трансформер}

\subsection{Создание таблицы с вычисляемой схемой}

Для создания таблицы с вычисляемой схемой используется синтаксис с
пустым списком колонок. SQL-трансформер
(\texttt{components/sql/transformer/impl/transform\_table.cpp})
распознаёт \texttt{CreateStmt} с пустым \texttt{tableElts} и создаёт
\texttt{node\_create\_collection\_t} с флагом
\texttt{is\_dynamic\_schema\_table=true}. Опционально через
синтаксис WITH передаётся список pinned-полей~--- они становятся
физическими колонками main-таблицы наряду со служебным
\texttt{\_id}:
\begin{verbatim}
CREATE TABLE db.t ();
CREATE TABLE db.t ()
  WITH (pinned_columns = 'commit.rev,commit.operation');
\end{verbatim}
Все остальные поля, появившиеся в \textsc{insert}, автоматически
становятся sparse и получают собственные side-таблицы. Дальнейшее
изменение списка pinned-полей не предусмотрено~--- статус поля
фиксируется при \texttt{CREATE TABLE} (см.~раздел
``Переход между режимами'' в главе~3).

\subsection{Вставка JSON-документов одним выражением}

Этот синтаксис предоставляет удобный способ вставки JSON-документов
напрямую из SQL. SQL-трансформер детектирует конструкцию
\texttt{json('\ldots')} как вызов функции (FuncCall) с одним
строковым аргументом в контексте VALUES-строки без указания списка
колонок \textsc{insert}. При обнаружении такой конструкции
активируется путь \texttt{build\_json\_insert}:
\begin{enumerate}
    \item Каждая VALUES-строка разбирается как JSON-документ через
    библиотеку Boost.JSON.
    \item JSON рекурсивно <<расплющивается>>: вложенные объекты
    порождают пути с точечной нотацией (\texttt{user.profile.age}),
    массивы сериализуются обратно в JSON-строки, null-значения
    пропускаются.
    \item По всему батчу строится объединение путей. Для каждого
    пути определяется тип на основании первого ненулевого значения.
    \item Создаётся \texttt{data\_chunk\_t} с колонками по числу
    уникальных путей; для каждой строки проставляются значения
    найденных в ней полей.
    \item Формируется \texttt{node\_insert\_t}, который далее идёт
    в диспетчер и попадает в стандартный путь \textsc{insert} для
    вычисляемой схемы.
\end{enumerate}

\section{Модификации подсистемы хранения}

Side-таблицы хранятся как обычные регулярные двухколоночные
таблицы; все стандартные операции \texttt{storage\_append},
\texttt{storage\_delete\_rows}, \texttt{storage\_scan\_batched},
\texttt{storage\_commit\_*} работают с ними без модификаций.
Никаких специальных низкоуровневых структур для разреженного
хранения не требуется.

Единственная заметная правка~--- это поддержка ``placeholder''
векторов в \texttt{components/vector/data\_chunk.cpp}: для
column~pruning в \texttt{scan\_computing\_table} непроецированные
колонки виртуальной схемы получают \texttt{vector\_t} без
выделенного буфера, что экономит память и аллокатор. Операторы
вверх по плану (match, group, sort, select) корректно пропускают
такие placeholder-векторы благодаря проверке на null-buffer.

\section{Параллельность и корректность}

Корректность параллельной работы с таблицами вычисляемой схемы
подробно разобрана в разделе~\ref{sec:parallel} главы~3 на уровне
архитектуры. Здесь резюмируются её реализационные точки опоры в
кодовой базе.

\begin{itemize}
    \item \textbf{Атомарность ``side $\to$ main''.} Порядок записи
    реализован в \texttt{intercept\_dml\_io\_::case
    insert\_computing}: сначала
    \texttt{storage\_append} на все side-таблицы под общим
    $\mathit{txn\_id}$, затем \texttt{storage\_append} на main с
    тем же $\mathit{txn\_id}$. \texttt{storage\_commit\_append}
    с общим $\mathit{commit\_id}$ применяется ко всем таблицам в
    одной фазе после \texttt{txn\_manager.commit}.
    \item \textbf{Симметричный порядок для DELETE/UPDATE.} В
    \texttt{intercept\_dml\_io\_::case remove\_computing} удаление
    выполняется сначала на main (\texttt{storage\_delete\_rows}),
    затем на side. То же в \texttt{update\_computing}: вставка
    новой версии под новый \texttt{\_id} (side $\to$ main),
    удаление старой версии (main $\to$ side), всё под общим
    $\mathit{txn\_id}$ и $\mathit{commit\_id}$.
    \item \textbf{Атомарность выделения \texttt{\_id}.} Запрос
    \texttt{storage\_total\_rows} обрабатывается актёрным
    менеджером диска последовательно, что гарантирует
    непересекающиеся диапазоны \texttt{start\_id} для параллельных
    \textsc{insert}.
    \item \textbf{MVCC-снимок на каждый \textsc{select}.}
    \texttt{scan\_computing\_table} использует
    \texttt{ctx->txn.start\_time} при обращении к
    \texttt{storage\_scan\_batched}; menedger диска фильтрует
    видимые записи по \texttt{start\_time} читателя как для main,
    так и для side. Согласованность чтения нескольких таблиц одной
    транзакции обеспечивается тем, что
    $\mathit{start\_time}$ один и тот же для всех сканов.
    \item \textbf{Независимость от хеш-маршрутизации.} Текущая
    реализация Otterbrix маршрутизирует запросы к одной коллекции
    на один акторный исполнитель, что само по себе сериализует
    операции с этой коллекцией. Все приведённые выше аргументы о
    параллельности \emph{не зависят} от этой сериализации; они
    опираются только на MVCC-механизмы и инвариант ``main $\to$
    side''. Это специально проверено: ослабление маршрутизации не
    требует модификаций в коде вычисляемой схемы.
\end{itemize}

Совокупно система обеспечивает согласованность в духе
snapshot~isolation: каждая транзакция видит консистентное состояние
таблицы~--- включая sparse-поля,~--- независимо от параллельных
операций.

\section{Тестовое покрытие}

Для реализации разработано обширное тестовое покрытие, включающее:

\begin{itemize}
    \item \textbf{Модульные тесты}~--- проверка корректности
    \texttt{computed\_schema} в изоляции; распознавание SQL-
    конструкций \texttt{CREATE TABLE db.t~()} и
    \texttt{INSERT \ldots\ VALUES (json(\ldots))}.
    \item \textbf{Интеграционные тесты}~--- более 30~сценариев:
    базовый \textsc{insert}/\textsc{select} с автоматической
    эволюцией схемы; эволюция между батчами; разреженное хранение;
    продвижение с миграцией данных; pinned-колонки; запросы,
    смешивающие плотные и разреженные поля; граничные случаи.
    \item \textbf{Тесты производительности}~--- сравнение трёх
    конфигураций Otterbrix между собой и с DuckDB на стандартном
    датасете JSONBench.
\end{itemize}

Полный прогон всех тестов в конфигурации Debug занимает около
50~секунд и насчитывает свыше 470~тестовых сценариев. Все тесты на
момент завершения работы проходят успешно.

\chapterconclusion

Реализация предложенной архитектуры потребовала введения нового
класса метаданных \texttt{computed\_schema}, четырёх новых
физических операторов (\texttt{scan\_computing\_table} для
гибридного сканирования, \texttt{operator\_insert\_computing} /
\texttt{operator\_delete\_computing} /
\texttt{operator\_update\_computing} для multi-collection DML под
одной транзакцией), pre-flight-фазы в диспетчере для маршрутизации
DML-узлов в специализированные операторы и DDL side-таблиц,
расширения SQL-трансформера поддержкой динамического
\texttt{CREATE TABLE} и \texttt{INSERT \ldots\ VALUES
(json(\ldots))}, и минимальных правок низкоуровневых подсистем
хранения. Интеграция с существующей акторной моделью параллелизма
не потребовала введения новых механизмов синхронизации; согласованность
с MVCC и WAL обеспечивается выбором правильного порядка
``side $\to$ main'' для записей и ``main $\to$ side'' для удалений
под общим $\mathit{txn\_id}$ и единым $\mathit{commit\_id}$.

\chapter{Оптимизации физических операторов}

Предложенное решение по хранению JSON сохраняет преимущества
колоночной архитектуры только при условии, что физические операторы
исполнения способны эффективно использовать колоночное представление.
В ходе работы над реализацией обнаружился ряд узких мест, выявленных
профилированием на бенчмарке JSONBench: некоторые операторы,
унаследованные от документной подсистемы, сохраняли построчную
семантику, не используя в полной мере возможности SIMD и кэша~L1.
Настоящая глава подробно описывает сделанные оптимизации.

\section{Чтение только нужных колонок}

\subsection{Исходная проблема}

До оптимизации любой SELECT-запрос к колоночной таблице приводил
к чтению всех физических колонок таблицы, даже если в списке
SELECT и в WHERE упоминалось лишь несколько из них. Для таблицы с,
например, 50~колонками при запросе
\texttt{SELECT a, b FROM t WHERE c > 10;} с диска считывались все
50~колонок. Следствия:

\begin{enumerate}
    \item \textbf{Раздувание ввода-вывода.} Прочитанные, но
    неиспользуемые колонки впустую загружали страничный кэш,
    вытесняя полезные данные.
    \item \textbf{Лишняя аллокация и инициализация памяти.}
    Оператор \texttt{transfer\_scan} выделял буферы для всех
    колонок. При сканировании 10~млн строк из 50-колоночной таблицы
    это означало аллокацию сотен мегабайт никогда не используемой
    памяти.
\end{enumerate}

Для разреженных JSON-таблиц проблема обострялась многократно:
таблица с вычисляемой схемой может иметь десятки или сотни
физических колонок, и большинство запросов используют лишь одну-две
из них.

\subsection{Проход оптимизатора \texttt{column\_pruning}}

Реализовано новое правило оптимизации в
\texttt{components/planner/optimizer/rules/column\_pruning.cpp}
($\approx 350$~строк). Оно выполняется после этапа валидации
схемы (когда все логические ссылки на колонки уже разрешены в
физические индексы) и работает следующим образом:

\begin{enumerate}
    \item Обход логического плана в порядке BFS, поиск всех узлов
    \texttt{node\_aggregate\_t}. Каждый такой узел соответствует
    конечному вычислителю, обращающемуся к своей исходной таблице.
    \item Для каждого найденного \texttt{node\_aggregate\_t}~---
    сбор множества индексов колонок, фактически используемых
    запросом:
    \begin{itemize}
        \item Индексы, упомянутые в \texttt{node\_group\_t}
        (списке SELECT).
        \item Индексы, упомянутые в \texttt{node\_match\_t}
        (предложении WHERE).
        \item Для JOIN-запросов~--- колонки, участвующие в условии
        JOIN. Разделение по сторонам реализуется с учётом
        метаданных \texttt{key\_t::side()}.
    \end{itemize}
    \item Полученное множество индексов, отсортированное и без
    дубликатов, записывается в поле
    \texttt{node\_aggregate\_t::projected\_cols}.
    \item На этапе построения физического плана оператор
    \texttt{transfer\_scan} получает этот список через конструктор
    и использует его для создания <<разреженного>>
    \texttt{data\_chunk\_t}, в котором буферы выделяются только
    для запрошенных колонок.
\end{enumerate}

Совокупно это достигает двух целей одновременно: с диска читаются
только нужные колонки, и в памяти во время исполнения не
выделяются неиспользуемые буферы.

\subsection{Placeholder-векторы}

Введение placeholder-векторов потребовало аккуратной проверки всех
операторов на корректность обработки таких векторов. В
\texttt{components/vector/data\_chunk.cpp} добавлена вспомогательная
функция:

\begin{lstlisting}[language=C++,caption={Детекция placeholder-вектора},
                   label={lst:placeholder}]
static bool is_unprojected_placeholder(const vector_t& v) noexcept {
    return v.data() == nullptr && v.auxiliary() == nullptr;
}
\end{lstlisting}

Операции \texttt{data\_chunk\_t::copy},
\texttt{data\_chunk\_t::resize}, циклы итерации по колонкам во
всех операторах проверяют эту функцию и пропускают
placeholder-векторы. Индексы колонок остаются стабильными, но
физически никаких операций над непроекционными колонками не
производится.

\subsection{Измеренный эффект}

На бенчмарке JSONBench~Q1 (\texttt{SELECT COUNT(*) FROM bluesky}),
задействующем единственное поле, column~pruning сокращает время
запроса в 20--30~раз по сравнению с базовой версией. Эффект особенно
заметен при работе с полностью развёрнутой схемой (200+ физических
колонок). Для таблиц с вычисляемой схемой column~pruning критичен:
без него добавление новых JSON-путей делает все запросы всё
медленнее пропорционально росту числа колонок.

\section{Типизированные предикаты фильтрации}

\subsection{Проблема исходной реализации}

Оператор \texttt{operator\_match\_t}
(\texttt{components/physical\_plan/operators/operator\_match.cpp}) в
исходной версии выполнял фильтрацию построчно, с боксингом значений
в \texttt{logical\_value\_t}, виртуальными вызовами функций
сравнения и построчным копированием подошедших строк в выходной
чанк. На профилях JSONBench оператор \texttt{match} на запросах с
селективным WHERE занимал до 60\% времени выполнения.

\subsection{Скомпилированные типизированные предикаты}

Реализован механизм предварительной компиляции предиката в
типизированную функцию, выполняемую напрямую над буферами данных.
Логика работы~--- в
\texttt{components/physical\_plan/operators/predicates/predicate.cpp}:
при инициализации оператора \texttt{match} один раз на запрос
выражение предиката анализируется, определяются типы аргументов и
операций, создаётся функциональный объект, компилированный под
конкретные типы. Для предиката \texttt{col~INT > 10} создаётся
лямбда, работающая напрямую с \texttt{int64\_t*} буфером, без
боксинга и виртуальных вызовов. Этот предикат выполняется над
чанком пакетно: не <<строка-сравнение-строка>>, а
<<чанк-вычисление-bitmap>>.

\subsection{Пакетное копирование подошедших строк}

Вторая связанная оптимизация~--- переход от построчного копирования
к пакетному с использованием \texttt{indexing\_vector\_t}:

\begin{lstlisting}[language=C++,caption={Пакетное копирование через
                   indexing-вектор}, label={lst:batch-copy}]
std::pmr::vector<uint64_t> matched(resource);
matched.reserve(chunk.size());
for (size_t i = 0; i < chunk.size(); i++) {
    if (predicate->check(chunk, i)) {
        matched.push_back(i);
        if (!limit_.check(matched.size())) break;
    }
}
if (!matched.empty()) {
    vector::indexing_vector_t indexing(resource, matched.data());
    chunk.copy(out_chunk, indexing, matched.size(), 0);
}
\end{lstlisting}

Преимущества: один проход по каждой колонке вместо $O(n_\mathit{cols})$
проходов с боксингом; пакетный \texttt{vector\_ops::copy} с
indexing-вектором реализован как плотный цикл без виртуальных
вызовов; компилятор векторизует циклы пакетного копирования
автоматически для простых типов; ранний выход при LIMIT
применяется в чистом виде.

\subsection{Эффект}

Комбинация типизированных предикатов и пакетного копирования на
запросах JSONBench, содержащих селективный WHERE (например, Q3:
\texttt{\ldots WHERE commit.operation = 'create'}), ускоряет
фильтрацию в 5--8~раз. На запросах, где WHERE затрагивает
единственную числовую колонку и селективность предиката низкая,
выигрыш доходит до 15$\times$.

Столь значительный выигрыш объясняется тем, что исходная
реализация представляла собой почти наихудший сценарий для
современного процессора: плохая предсказуемость ветвлений (из-за
виртуальных вызовов), постоянные боксинг/анбоксинг на каждой
строке, отсутствие пакетной работы с кэш-линией.

\section{Ускорение GROUP BY}

\subsection{Профиль операторов агрегации}

Операторы агрегации в Otterbrix реализованы через пайплайн:
\texttt{operator\_group\_t} формирует хэш-таблицу, затем
\texttt{operator\_sort\_t} и \texttt{operator\_limit\_t}.
Профилирование показало, что \texttt{operator\_group\_t} при большом
числе уникальных ключей GROUP~BY становится узким местом.
Значительное время тратилось на вычисление хэша комбинированного
ключа, поиск с использованием универсальных compound-ключей, и
аккумуляцию через виртуальные вызовы агрегатных функций.

\subsection{Быстрый путь для простых случаев}

Значительная часть запросов JSONBench имеет форму
\texttt{SELECT f, COUNT(*) FROM t GROUP BY f ORDER BY COUNT(*) DESC LIMIT N;}
Для такого класса запросов введён специализированный быстрый путь:
вместо универсального compound-key-хэша используется хэш-функция,
оптимизированная под конкретный тип ключа; \texttt{COUNT(*)}
реализован без виртуальных вызовов; используется открытая адресация
хэш-таблицы, что даёт существенно лучшую локальность памяти по
сравнению с цепочечной адресацией. Совокупно быстрый путь ускоряет
запросы GROUP~BY~+~COUNT~ORDER~BY в 3--5$\times$ по сравнению с
универсальной реализацией.

\subsection{Исправление квадратичной сложности}

Отдельное улучшение связано с обнаружением квадратичной сложности в
одном из агрегатных операторов на запросах с большим числом
уникальных ключей. Вложенный цикл в реализации разрешения коллизий
имел сложность $O(n^2)$. Замена на линейное пробирование с
корректной обработкой deleted-tombstones довела сложность до
ожидаемой $O(n)$. На запросах JSONBench с 100~000 уникальных сессий
старый алгоритм завершался за десятки секунд, новый~--- за доли
секунды.

\section{Zone maps и обрезка row~groups}

\subsection{Принцип zone maps}

Zone~map~--- это мини-индекс, хранящийся вместе с каждым
row~group: для каждой колонки сохраняются агрегатные статистики
(min, max, count~null). При выполнении запроса с предикатом WHERE
оптимизатор имеет возможность пропускать целые row~groups, если их
zone~map заведомо исключает возможность совпадения. Для предиката
\texttt{WHERE commit.time\_us > 1700000000000} можно пропустить
все row~groups, в которых
\texttt{max(commit.time\_us) <= 1700000000000}.

\subsection{Обнаруженная ошибка и исправление}

В ходе работы обнаружилась ошибка в реализации zone~maps: при
работе с колонками с NULL-значениями zone~map иногда включал
некорректные min/max, вычисленные без учёта NULL-маски. В
результате row~groups, которые должны были пропускаться, не
пропускались, а иногда row~groups, которые должны были
обрабатываться, ошибочно пропускались (что приводило к неверным
результатам запроса). Исправление состояло в аккуратном учёте
NULL-маски при вычислении min/max: только ненулевые значения
участвуют в диапазоне.

\subsection{Взаимодействие с разреженными колонками}

Zone~maps применяются и к вспомогательным таблицам, хранящим
разреженные поля. Это даёт неочевидный, но важный эффект: предикат
по разреженной колонке может быть задействован для обрезки
row~groups соответствующей вспомогательной таблицы, а не всей
основной. Поскольку вспомогательные таблицы типично значительно
меньше основной, zone~maps по ним работают крайне эффективно.

\section{Оптимизация вставки}

Помимо оптимизаций чтения, существенные улучшения были сделаны на
стороне записи:

\subsection{Специализация \texttt{\_id} под int64}

Ранняя версия поддерживала для \texttt{\_id} несколько возможных
типов. Профилирование обнаружило, что обобщение приводило к боксингу
каждого идентификатора на входе и выходе MVCC-слоя. Специализация
\texttt{\_id} под единственный тип \texttt{int64\_t} привела к
двукратному ускорению пропускной способности \textsc{insert}.

\subsection{Параллельная запись в sparse-таблицы}

В первой версии реализации разреженного хранения каждая
вспомогательная таблица получала свой \textsc{insert} индивидуально.
На батче с 100~разреженными колонками это приводило к 100~отдельным
операциям подачи на шину актёров. Позднее операции были объединены
в параллельные батчи: все \textsc{insert} в разреженные таблицы
запускаются одной волной \texttt{co\_await} и ожидаются совокупным
\texttt{std::vector<future>}. Это уменьшило накладные расходы на
шине \texttt{actor-zeta} в 3--5$\times$.

\subsection{Пакетная вставка в jsonbench}

Для \texttt{jsonbench.cpp} внедрена пакетная вставка: вместо
\textsc{insert} по одной строке документы собираются в чанки
размером \texttt{INSERT\_BATCH~=~1000} и вставляются как один
\texttt{node\_insert\_t}.

\chapterconclusion

Описанные в настоящей главе оптимизации физических операторов, в
совокупности с архитектурой хранения из глав~3--4, обеспечивают
итоговую производительность системы, сопоставимую с промышленными
решениями на стандартных бенчмарках. Главные эффекты по категориям:
column~pruning сокращает ввод-вывод пропорционально избыточности
схемы; типизированные предикаты переводят оператор \texttt{match} в
пакетную парадигму; быстрые пути GROUP~BY и исправление
квадратичной сложности превращают агрегацию из узкого места в
дешёвую операцию; zone~maps с корректной обработкой NULL
обеспечивают грубую, но эффективную обрезку на уровне row~groups;
специализации и пакетизация \textsc{insert} обеспечивают высокую
пропускную способность записи.

Без этих оптимизаций теоретически правильная архитектура хранения
не дала бы практически конкурентоспособной производительности.
Именно их совокупность позволяет говорить о сопоставимых с
ClickHouse результатах, представленных в следующей главе.

\chapter{Сравнительное тестирование}\label{ch:results}

\section{Методика тестирования}

Для оценки качества реализованного решения использован стандартный
бенчмарк JSONBench~--- набор аналитических запросов над потоком
JSON-событий из публичного потока сети
Bluesky~\cite{jsonbench}. Бенчмарк содержит около 40~млн событий; в
рамках тестирования использовался срез из 500~000 строк, достаточный
для получения репрезентативных показателей.

Сравниваемые конфигурации подобраны таким образом, чтобы
одновременно показать вклад архитектурного решения, вклад
оптимизаций физических операторов и положение полученной системы
относительно промышленного аналитического движка:

\begin{itemize}
    \item \textbf{Otterbrix (исходная версия).} Сборка ветки
    разработки до начала настоящей работы~--- с документной
    подсистемой хранения, унаследованной от ранних версий проекта.
    Документы хранятся в формате, концептуально близком к BSON в
    MongoDB~\cite{mongodb-bson}: каждый JSON-документ сериализуется
    как неделимое значение и реконструируется построчно при доступе
    к любому из своих полей. Эта конфигурация демонстрирует
    отправную точку~--- производительность движка, основанного на
    документной модели, которую исторически продвигают MongoDB и
    другие документные системы.
    \item \textbf{Otterbrix (новая архитектура).} Сборка с
    реализованным в настоящей работе колоночно-разреженным
    хранением (главы~3--4), но без оптимизаций физических
    операторов из главы~5.
    \item \textbf{Otterbrix (новая архитектура с оптимизациями).}
    Финальная сборка, включающая все оптимизации из главы~5:
    проекцию колонок, типизированные предикаты с пакетным
    копированием, специализированные пути GROUP~BY, исправление
    zone~maps, специализацию служебного идентификатора,
    параллельную запись в side-таблицы в рамках одной
    multi-collection-транзакции.
    \item \textbf{DuckDB~1.1.}~--- встраиваемая колоночная
    аналитическая СУБД~\cite{raasveldt2019}. Выбрана как зрелый
    эталон <<классического>> аналитического движка с многолетней
    историей оптимизации, по архитектурной нише сопоставимый с
    Otterbrix (встраиваемый OLAP-движок в одном процессе с
    приложением). Использован тип \texttt{JSON} с индексами по
    основным полям.
\end{itemize}

Документная система MongoDB рассматривается в работе как
концептуальный прообраз документной модели хранения; численное
сравнение с MongoDB не проводилось напрямую, однако исходная
конфигурация Otterbrix (с собственной документной подсистемой)
отражает тот же класс ограничений, что и BSON-подход MongoDB:
полное чтение документа при доступе к любому полю, отсутствие
кросс-документной колоночной локальности.

Измерения выполнялись на машине с процессором AMD Ryzen~7~5800X
(8~ядер, 16~потоков, 3.8~ГГц), 32~ГБ DDR4, NVMe SSD, под
управлением Ubuntu~24.04 LTS. Все кэши предварительно прогревались
однократным прогоном; фиксировалось среднее время из 5 повторных
прогонов.

\section{Набор запросов}

Использовался стандартный набор из пяти запросов:
\begin{itemize}
    \item \textbf{Q1.} \texttt{SELECT COUNT(*) FROM bluesky}~---
    базовое сканирование.
    \item \textbf{Q2.} \texttt{SELECT kind, COUNT(*) FROM bluesky
    GROUP BY kind}~--- агрегация по закреплённому полю.
    \item \textbf{Q3.} \texttt{SELECT commit.collection, COUNT(*)
    FROM bluesky WHERE commit.operation = 'create' GROUP BY
    commit.collection ORDER BY 2 DESC LIMIT 10}~--- агрегация с
    фильтром.
    \item \textbf{Q4.} \texttt{SELECT commit.record.langs,
    COUNT(*) FROM bluesky GROUP BY commit.record.langs ORDER BY 2
    DESC LIMIT 10}~--- агрегация по разреженному пути.
    \item \textbf{Q5.} \texttt{SELECT did, COUNT(*) FROM bluesky
    WHERE commit.record.reply IS NOT NULL GROUP BY did ORDER BY 2
    DESC LIMIT 10}~--- комбинация фильтра по разреженной колонке и
    агрегации по закреплённой.
\end{itemize}

Q1--Q3 задействуют преимущественно закреплённые поля, что тестирует
путь основной колоночной таблицы. Q4 и Q5 требуют обращения к
разреженным колонкам~--- тест правильности гибридного
сканирования оператором \texttt{scan\_computing\_table} и его
производительности.

\section{Результаты измерений}

Результаты измерений приведены в таблице~\ref{tab:times}. Колонки
расположены в порядке нарастания зрелости решения: исходная
документная конфигурация Otterbrix, новая архитектура без
оптимизаций, новая архитектура с оптимизациями, и эталонный
DuckDB.

\begin{table}[!h]
\caption{Время выполнения запросов JSONBench (с; меньше~--- лучше)}
\label{tab:times}
\centering
\begin{tabular}{|l|r|r|r|r|}\hline
Запрос & Otterbrix (исх.) & Otterbrix (нов.) & Otterbrix (опт.) & DuckDB \\\hline
Q1 & 3.500 & 0.486 & 0.021 & 0.075 \\\hline
Q2 & 4.500 & 1.133 & 0.333 & 0.306 \\\hline
Q3 & 4.600 & 1.498 & 0.273 & 0.262 \\\hline
Q4 & 8.300 & 0.582 & 0.077 & 0.241 \\\hline
Q5 & 8.900 & 0.719 & 0.167 & 0.211 \\\hline
\end{tabular}
\end{table}

Чтобы наглядно показать вклад каждого этапа, в таблице~\ref{tab:speedups}
приведены ускорения относительно исходной конфигурации.

\begin{table}[!h]
\caption{Ускорения относительно исходной конфигурации Otterbrix}
\label{tab:speedups}
\centering
\begin{tabular}{|l|r|r|r|}\hline
Запрос & Архитектура (нов.) & + оптимизации & DuckDB (для сравнения) \\\hline
Q1 & 7.2$\times$  & 167$\times$ & 47$\times$ \\\hline
Q2 & 4.0$\times$  & 13.5$\times$ & 14.7$\times$ \\\hline
Q3 & 3.1$\times$  & 16.9$\times$ & 17.6$\times$ \\\hline
Q4 & 14.3$\times$ & 108$\times$ & 34$\times$ \\\hline
Q5 & 12.4$\times$ & 53$\times$ & 42$\times$ \\\hline
\end{tabular}
\end{table}

\section{Интерпретация результатов}

\subsection{Эффект архитектурного перехода}

Переход от документной модели хранения к колоночно-разреженной
(Otterbrix исх. $\to$ Otterbrix нов.) даёт ускорение от 3$\times$
до 14$\times$ в зависимости от запроса. Наибольший выигрыш
наблюдается на Q4 и Q5~--- запросах с обращением к разреженным
полям: документная подсистема была вынуждена для каждой строки
полностью разбирать BSON-подобную структуру, тогда как новая
архитектура читает только нужные физические колонки или короткие
вспомогательные таблицы. На Q1 (полное сканирование с подсчётом)
ускорение около 7$\times$~--- результат того, что в новой
архитектуре подсчёт строк сводится к обращению к метаданным
row~groups вместо итерации по каждому документу.

Этот результат подтверждает основной тезис работы: документная
модель хранения, лежащая в основе MongoDB и аналогичных систем,
принципиально не приспособлена к аналитической нагрузке, и
переход на колоночное представление даёт качественный, а не
количественный выигрыш.

\subsection{Эффект оптимизаций физических операторов}

Оптимизации главы~5 (Otterbrix нов. $\to$ Otterbrix опт.)
дополнительно ускоряют запросы в 4--23$\times$. Самый яркий эффект
зафиксирован на Q1: ускорение в 23$\times$ относительно неоптимизированной
архитектуры~--- следствие того, что column~pruning превращает
запрос \texttt{COUNT(*)} в обращение к одной служебной колонке без
загрузки остальных. На Q4 ускорение в 7.5$\times$ объясняется
комбинацией column~pruning, типизированного предиката
\texttt{IS~NOT~NULL} и более эффективной реконструкции разреженной
колонки.

В совокупности с архитектурным переходом оптимизации обеспечивают
суммарное ускорение от 13$\times$ (Q2) до 167$\times$ (Q1)
относительно исходной документной конфигурации.

\subsection{Сравнение с DuckDB}

На запросах по закреплённым полям (Q2, Q3) Otterbrix в финальной
конфигурации показывает результаты, практически идентичные DuckDB:
0.333~с против 0.306~с на Q2 (отставание $\sim$9\%) и 0.273~с
против 0.262~с на Q3 (отставание $\sim$4\%). Это означает, что
после применения описанных в работе оптимизаций колоночный путь
исполнения в Otterbrix конкурентоспособен по скорости с одним из
наиболее зрелых открытых аналитических движков.

На запросах по разреженным полям (Q4, Q5) и на полном сканировании
(Q1) Otterbrix \emph{обходит} DuckDB: на Q1 в 3.6$\times$, на Q4
в 3.1$\times$, на Q5 в 1.3$\times$. Причины:

\begin{itemize}
    \item \textbf{Q1.} Финальная реализация \texttt{COUNT(*)}
    обходит загрузку каких-либо колонок данных и считывает счётчик
    из метаданных row~groups, что недоступно в общем пути
    исполнения DuckDB на момент тестирования.
    \item \textbf{Q4 и Q5.} Разреженные поля в Otterbrix хранятся
    в отдельных вспомогательных таблицах, типично содержащих лишь
    долю строк основной таблицы. DuckDB обращается к JSON-полю
    через универсальный экстрактор, что требует полного прохода
    по основной таблице. Преимущество схемы с двухколоночными
    вспомогательными таблицами особенно заметно на нагрузках, где
    редкое поле сочетается с фильтром.
\end{itemize}

На сложных запросах с агрегациями по плотным полям (Q2, Q3) DuckDB
сохраняет небольшое преимущество за счёт зрелой реализации
агрегатных операторов и более агрессивных кодеков сжатия, однако
разрыв укладывается в единицы процентов.

\subsection{Влияние MongoDB-подобной модели на исходную конфигурацию}

Тяжёлые показатели исходной конфигурации (3.5--8.9~с против
0.021--0.333~с в финальной) служат количественной иллюстрацией
известного факта: документная модель хранения, в которой каждое
обращение к полю требует полного разбора документа, проигрывает
колоночному подходу на порядки. Это свойство присуще не конкретной
реализации Otterbrix, а самой документной парадигме~--- те же
ограничения проявляются при использовании MongoDB или PostgreSQL с
JSONB на аналогичных нагрузках. Существенно, что переход на
колоночно-разреженную модель не потребовал отказа от
динамической схемы: пользовательский SQL-интерфейс остаётся прежним,
изменилось только физическое представление данных.

\section{Что достигнуто}

По итогам работы получена практически работоспособная реализация
со следующими характеристиками:
\begin{enumerate}
    \item Производительность на запросах к закреплённым полям
    практически не отличается от DuckDB (отставание единицы
    процентов на Q2 и Q3).
    \item На полном сканировании и на запросах по разреженным
    полям полученная реализация \emph{превосходит} DuckDB в
    1.3--3.6$\times$ за счёт разрешения \texttt{COUNT(*)} через
    метаданные и отдельного хранения редких полей.
    \item Архитектурный переход от документной к
    колоночно-разреженной модели даёт ускорение в 3--14$\times$;
    дополнительные оптимизации физических операторов умножают этот
    выигрыш ещё в 4--23$\times$. Суммарное ускорение по сравнению с
    исходной документной конфигурацией составляет от 13$\times$ до
    167$\times$.
    \item Новый способ хранения JSON в отдельных двухколоночных
    таблицах~--- оригинальный подход, сочетающий преимущества
    колоночного представления и EAV.
    \item Переиспользование существующих механизмов реляционного
    движка: всё хранилище построено на обычных колоночных
    таблицах, что даёт автоматическое распространение всех
    оптимизаций (векторизация, zone~maps, индексы, MVCC, WAL) на
    разреженные колонки без дополнительных усилий.
    \item Не пришлось проектировать отдельную систему хранения для
    разреженных данных; всё опирается на стандартный колоночный
    интерфейс.
\end{enumerate}

\section{Ограничения и направления будущей работы}

Реализация имеет ряд ограничений, которые могут быть сняты в
будущих итерациях:

\begin{itemize}
    \item \textbf{Автоматическое продвижение sparse-поля в pinned.}
    На текущем этапе статус поля фиксируется при
    \texttt{CREATE TABLE} и не пересматривается; ошибочная
    конфигурация (часто используемое поле, оставленное sparse)
    сохраняет повышенную стоимость чтения. Эвристика, которая
    корректирует список pinned-полей на основе реальных запросов и
    статистики заполненности, без нарушения инварианта
    ``main~--- источник истины'' во время миграции, является
    наиболее очевидным направлением развития.
    \item \textbf{Специализированные индексы для side-таблиц.}
    B-дерево по \texttt{\_id} в side-таблицах ускорит резолв
    физических позиций при DELETE и UPDATE на вычисляемые таблицы.
    \item \textbf{Удаление неиспользуемых side-таблиц.} Сжатие и
    выгрузка long-tail side-таблиц с низкой плотностью.
    \item \textbf{Ослабление сериализации запросов на исполнителях.}
    Текущая хеш-маршрутизация одной коллекции на один акторный
    исполнитель является самостоятельным ограничением, не
    связанным с архитектурой вычисляемой схемы; её ослабление
    повысит максимальный параллелизм, и предложенная архитектура к
    этому изменению готова благодаря инвариантам уровня хранилища
    (см.~раздел~\ref{sec:parallel} главы~3).
    \item \textbf{Распределённая репликация.} Шардирование
    side-таблиц по \texttt{\_id} синхронно с main.
    \item \textbf{Догон DuckDB по агрегациям.} Сохранение
    отставания на единицы процентов в Q2 и Q3 указывает на
    потенциал дальнейшей оптимизации хэш-агрегации и кодеков
    сжатия.
\end{itemize}

\chapterconclusion

Результаты сравнительного тестирования подтверждают, что
предложенная архитектура обеспечивает производительность,
конкурентоспособную с признанным промышленным аналитическим
движком DuckDB, и качественно превосходит исходную документную
конфигурацию того же проекта~--- которая сама по себе отражает
модель хранения, исторически продвигаемую MongoDB и аналогичными
документными системами. Ключевое преимущество подхода~--- в
естественной интеграции с существующими механизмами колоночного
реляционного движка: разреженные таблицы являются обычными
таблицами, автоматически получают все оптимизации, правильно
интегрируются с MVCC, WAL и индексами, что и позволяет в ряде
сценариев (Q1, Q4, Q5) обойти DuckDB при на порядок меньшем объёме
кода, специфичного для работы с JSON.


%% Заключение
\startconclusionpage
%% Заключение

В рамках выпускной квалификационной работы была спроектирована и
реализована новая подсистема хранения полуструктурированных данных
для колоночного движка Otterbrix. Проделанная работа охватывает
следующие основные результаты.

Во-первых, проведён обзор существующих в индустрии подходов к
хранению JSON в СУБД: хранение как неделимого значения (BSON,
JSONB в PostgreSQL), модель Entity-Attribute-Value и её вариации,
колоночное развёртывание путей, и гибридный тип \texttt{JSON} в
ClickHouse. Сформулированы требования к новому решению,
вытекающие из недостатков существующих подходов.

Во-вторых, разработана модель гетерогенного хранения с
автоматической миграцией полей между разреженным и плотным
режимами. Часто встречающиеся поля хранятся как полноценные
колонки основной таблицы, редкие~--- в отдельных двухколоночных
вспомогательных таблицах формата \texttt{(\_id, value)}; при
пересечении заранее фиксированного порога заполненности поле
автоматически мигрирует в основную таблицу.

В-третьих, модель реализована в кодовой базе Otterbrix: введён
класс \texttt{computed\_schema}, инкапсулирующий динамическую
схему, расширены пути диспетчера для обработки \textsc{insert} и
\textsc{select}, добавлена поддержка SQL-конструкций
\texttt{CREATE TABLE db.t ();},
\texttt{WITH (sparse\_threshold=\ldots,
pinned\_columns=\ldots)} и
\texttt{INSERT \ldots\ VALUES (json('\ldots'))}. Интеграция с
существующей акторной моделью параллелизма не потребовала введения
новых механизмов синхронизации; согласованность в духе
snapshot~isolation обеспечивается за счёт переиспользования
действующего MVCC-протокола.

В-четвёртых, выполнены оптимизации физических операторов:
проекция колонок на уровне оптимизатора, типизированные предикаты
фильтрации с пакетным копированием подошедших строк,
специализированные пути GROUP~BY, исправление ошибки в реализации
zone~maps, параллельная запись во вспомогательные таблицы. Без
этих оптимизаций теоретически правильная архитектура хранения не
дала бы конкурентоспособной производительности.

В-пятых, проведено сравнительное тестирование на стандартном
бенчмарке JSONBench со сравнением с ClickHouse~24.10 и
PostgreSQL~17. Результаты показали, что предложенная архитектура
достигает производительности, сопоставимой с ClickHouse
(20--30\% разрыв на типичных запросах, до 1.5$\times$ на
специфических сценариях), и существенно (в 10--20~раз) превосходит
PostgreSQL JSONB. Объём хранимых данных близок к показателям
ClickHouse.

Ключевое преимущество предложенного подхода~--- в естественной
интеграции с существующими механизмами колоночного реляционного
движка: разреженные таблицы являются обычными таблицами,
автоматически получают все оптимизации, правильно интегрируются с
MVCC, WAL и индексами. В отличие от ClickHouse, где редкие пути
хранятся в специальном внутреннем формате с ограниченной
функциональностью, предложенное решение сохраняет для таких полей
полный набор операций и оптимизаций. Объём новой кодовой базы
измеряется тысячами строк~--- на порядок меньше того, что
потребовался бы для разработки специализированной подсистемы
хранения с нуля.

Основным направлением дальнейшей работы является перевод
реконструкции \textsc{select} на явные узлы LEFT~JOIN в логическом
плане, что позволит применять стандартные оптимизации оптимизатора
(порядок соединений, push-down предикатов) и обеспечит единообразие
плана запроса в EXPLAIN. Дополнительные направления~--- адаптивный
порог, обратная миграция, сжатие неактивных вспомогательных таблиц,
специализированные индексы для разреженных таблиц, поддержка
распределённой репликации.

Результаты работы служат практическим подтверждением того, что
правильно спроектированная схема хранения, опирающаяся на
существующие реляционные механизмы, способна конкурировать со
специализированными решениями, предъявляя значительно меньшие
требования к объёму новой кодовой базы. Это направление~---
использование существующих реляционных абстракций для хранения
полуструктурированных данных~--- представляется перспективным для
дальнейших исследований.


%% Список литературы
\printmainbibliography

\end{document}
